\documentclass[11pt,a4paper]{article}

\usepackage[utf8]{inputenc}
\usepackage[T1]{fontenc}
\usepackage{amsmath,amssymb,amsthm,amsfonts}
\usepackage{bm}
\usepackage[margin=1in]{geometry}
\usepackage{url}

\newcommand{\toprule}{\hline}
\newcommand{\midrule}{\hline}
\newcommand{\bottomrule}{\hline}

\newtheorem{theorem}{Theorem}[section]
\newtheorem{lemma}[theorem]{Lemma}
\newtheorem{proposition}[theorem]{Proposition}
\newtheorem{corollary}[theorem]{Corollary}
\newtheorem{definition}[theorem]{Definition}
\newtheorem{remark}[theorem]{Remark}
\newtheorem{assumption}{Assumption}
\newtheorem{objection}{Objection}
\newtheorem{response}{Response}

\DeclareMathOperator*{\argmax}{arg\,max}
\DeclareMathOperator*{\argmin}{arg\,min}
\DeclareMathOperator{\Tr}{Tr}
\DeclareMathOperator{\KL}{KL}
\DeclareMathOperator{\Var}{Var}
\DeclareMathOperator{\Cov}{Cov}
\DeclareMathOperator{\supp}{supp}
\DeclareMathOperator{\rank}{rank}
\DeclareMathOperator{\diam}{diam}

\title{\textbf{On the Impossibility of Autonomous Self-Improvement\\to General Intelligence in Stochastic Language Models:\\A Proof Program via Information Theory, G\"{o}delian Limits,\\and Distributional Collapse}\\[10pt]\large Revised Edition with Counterargument Analysis and Empirical Audit}
\author{Filip Pajalic}
\date{Draft v6, August 2026}

\begin{document}
\maketitle

\begin{abstract}
We develop a proof program for the claim that large language models (LLMs) based on stochastic next-token prediction cannot, through autonomous self-improvement in a sealed loop, converge to artificial general intelligence (AGI). The model is treated as a conditional probability estimator trained by maximum likelihood, and the self-improvement loop as a composition of training channels in an explicit Bayesian joint in which ground truth is a random variable. In that setting we prove an information ceiling, convert it into an \emph{accuracy} ceiling by a Fano argument, and derive the impossibility of convergence for any starting model short of the target; this chain is machine-checked in Lean~4 with no \texttt{sorry} and no project-local axioms. Around this core we state seven further barriers---effective-support and entropy contraction, G\"{o}delian and Rice-theoretic bounds on self-verification, the self-distillation fixed point, the pattern-completion and complexity barriers, resource bounds on in-context self-modification, and non-improvement of the error recurrence---each labeled with its formalization status.

This version substantially narrows several earlier claims in response to adversarial review. The AGI criterion is relativized to a fixed task family, because the unrestricted ``matches every computable function'' criterion is unsatisfiable by diagonalization and would make the main theorem vacuous. Support contraction is restated over $\epsilon$-effective support, since literal support inclusion is trivial for softmax models. Error divergence is replaced by non-improvement as the universal claim, with divergence confined to the replacement regime; in the data-accumulation regime error is provably bounded. The context-window results are rescoped as bounds on in-context self-modification rather than on the loop in general. The load-bearing concept is reformulated from ``human influence'' to \emph{external grounding}, with human design as the dominant special case, and the boundary cases this concedes---a proof-checking kernel, an embodied agent---are stated explicitly. A concluding empirical section audits the 2023--2026 record in both directions, including results in tension with the claims, and states in advance what would falsify the program.
\end{abstract}

\tableofcontents
\newpage

%% ============================================================
\section{Scope and Definitions}
%% ============================================================

We begin by precisely delineating what this proof does and does not claim.

\begin{remark}[Formalization Status]
This manuscript is accompanied by a Lean~4 formalization. As of v6 the conditional core is machine-checked: the loop is defined as a composition of training channels, and the information ceiling, the entropy and effective-support contractions, the accuracy link, the error-recurrence results, and the main theorem are proved with no \texttt{sorry}. An axiom audit (\texttt{lean4/Impossibility/Audit.lean}) reports that these results depend only on Lean's own axioms (\texttt{propext}, \texttt{Classical.choice}, \texttt{Quot.sound}). Three named bridge axioms remain, none of them used by the main theorem: Kolmogorov complexity and its unboundedness, G\"{o}del's second incompleteness theorem over an abstract formal-system interface, and the modeling claim that a finite-precision LLM is a program. What is proved is an implication: the barriers follow \emph{from} the channel model of the loop, whose data-processing property is carried as an explicit hypothesis rather than derived from a measure-theoretic construction of the training process. The claims below should therefore be read in three categories:
\begin{enumerate}
    \item \emph{Formalized components}, such as selected computability facts imported from Mathlib;
    \item \emph{Conditional theorem targets}, whose conclusions depend on assumptions that must still be discharged;
    \item \emph{Empirical or philosophical support}, which may motivate the theorem targets but does not substitute for formal proof.
\end{enumerate}
The verification plan for the remaining gaps is kept in the repository file \texttt{docs/\allowbreak verification\_\allowbreak plan.md}. Two of the eight barriers---pattern completion and computational complexity---are deliberately absent from the Lean barrier structure: neither Kolmogorov complexity nor $\mathrm{P} \neq \mathrm{NP}$ has a Lean~4 formalization, and a contentless placeholder would misrepresent their status.
\end{remark}

\begin{remark}[Scope]
This proof program addresses the following specific question: \emph{Can a system based on autoregressive stochastic language modeling, trained via maximum likelihood estimation (or equivalent objectives), autonomously self-improve to AGI without any form of human influence at any stage of the self-improvement process?}

We do \textbf{not} claim:
\begin{itemize}
    \item That AGI is impossible by any means whatsoever.
    \item That hybrid systems combining LLMs with other components cannot advance.
    \item That human-guided improvement of LLMs is bounded.
\end{itemize}

We \textbf{do} target the following conditional claim: within the stochastic autoregressive framework operating autonomously, the mathematical structure of the system prevents convergence to AGI once the stated barriers are fully formalized. Any system that escapes these bounds should be analyzed by asking whether it incorporates elements that constitute human influence or departs from the stochastic language modeling paradigm.
\end{remark}

\subsection{The Large Language Model as a Stochastic System}

\begin{definition}[Autoregressive Language Model]
An autoregressive language model is a parameterized conditional probability distribution over a vocabulary $\mathcal{V} = \{v_1, v_2, \ldots, v_V\}$:
\begin{equation}
p_{\bm{\theta}}(x_t \mid x_{<t}) = \mathrm{softmax}\!\left(\bm{W}_{\mathrm{out}} \cdot h_L(\bm{x}_{<t}; \bm{\theta})\right)_{x_t}
\end{equation}
where $\bm{\theta} \in \mathbb{R}^p$ is the parameter vector (with $p$ on the order of $10^{9}$ to $10^{12}$), and $h_L$ is the output of an $L$-layer transformer network. The full sequence probability factorizes as:
\begin{equation}
p_{\bm{\theta}}(\bm{x}) = \prod_{t=1}^{T} p_{\bm{\theta}}(x_t \mid x_{<t})
\end{equation}
\end{definition}

\begin{definition}[Transformer Computation]
Each transformer layer $\ell \in \{1, \ldots, L\}$ computes:
\begin{align}
\bm{z}^{(\ell)} &= \mathrm{LayerNorm}\!\left(\bm{h}^{(\ell-1)} + \mathrm{MultiHeadAttn}(\bm{h}^{(\ell-1)})\right) \\
\bm{h}^{(\ell)} &= \mathrm{LayerNorm}\!\left(\bm{z}^{(\ell)} + \mathrm{FFN}(\bm{z}^{(\ell)})\right)
\end{align}
where multi-head attention with $H$ heads is:
\begin{equation}
\mathrm{MultiHeadAttn}(\bm{h}) = \bm{W}_O \!\left[\bigoplus_{i=1}^{H} \mathrm{softmax}\!\left(\frac{\bm{Q}_i \bm{K}_i^\top}{\sqrt{d_k}}\right) \bm{V}_i \right]
\end{equation}
with $\bm{Q}_i = \bm{W}_Q^{(i)} \bm{h}$, $\bm{K}_i = \bm{W}_K^{(i)} \bm{h}$, $\bm{V}_i = \bm{W}_V^{(i)} \bm{h}$, and:
\begin{equation}
\mathrm{FFN}(\bm{z}) = \bm{W}_2 \cdot \sigma(\bm{W}_1 \bm{z} + \bm{b}_1) + \bm{b}_2
\end{equation}
for activation function $\sigma$ (typically GELU or SiLU).
\end{definition}

\begin{definition}[Training Objective]
The model is trained to minimize the cross-entropy loss:
\begin{equation}
\mathcal{L}(\bm{\theta}) = -\mathbb{E}_{\bm{x} \sim \mathcal{D}_{\mathrm{train}}}\!\left[\frac{1}{T}\sum_{t=1}^{T} \log p_{\bm{\theta}}(x_t \mid x_{<t})\right] = H(p_{\mathrm{data}}) + \KL(p_{\mathrm{data}} \| p_{\bm{\theta}})
\label{eq:loss}
\end{equation}
\end{definition}

\subsection{Defining Autonomous Self-Improvement}

\begin{definition}[External Grounding Signal]
\label{def:grounding}
Work in the Bayesian joint of Section~\ref{sec:joint}, in which ground truth is a random variable $\Theta$. A signal $S$ available to the process at step $k \geq 1$ is an \emph{external grounding signal} (relative to $\bm{\theta}_k$) if
\begin{equation}
I(S; \Theta \mid \bm{\theta}_k) > 0,
\end{equation}
that is, if it carries information about ground truth that the model does not already contain. Signals with $I(S; \Theta \mid \bm{\theta}_k) = 0$---pseudo-random noise, re-samplings of the model's own distribution, permutations of its own outputs---are \emph{ungrounded}, however statistically independent of the model's errors they may be.
\end{definition}

\begin{definition}[Human Influence]
\label{def:human}
\emph{Human influence} is external grounding of human origin: an external grounding signal whose information about $\Theta$ originates in human cognition or in human-directed measurement. Concretely, any of the following introduced at step $k \geq 1$:
\begin{enumerate}
    \item New human-curated or human-generated training data,
    \item Human-designed reward functions, preference labels, or evaluation criteria,
    \item Human-designed filtering, selection, or ranking mechanisms applied to model outputs,
    \item Human-designed tools, environments, or external verifiers (including game rules, physics engines, code interpreters, proof checkers),
    \item Human-designed architectural modifications or hyperparameter schedules.
\end{enumerate}
Items (2)--(4) include systems where humans designed the evaluation mechanism \emph{before} the self-improvement loop begins. A reward function crafted by humans is human influence even if no human is present during training: the intelligence embedded in the design of the feedback signal originates from human cognition.
\end{definition}

\begin{remark}[Why grounding, not humanity, is load-bearing]
\label{rem:grounding}
Every formal argument in this manuscript distinguishes \emph{internal} from \emph{external} signal; none distinguishes \emph{human} from \emph{non-human} origin. The data processing inequality does not know who wrote the reward function. We therefore state the barriers in terms of Definition~\ref{def:grounding} and treat Definition~\ref{def:human} as the empirically dominant special case: in every self-improvement system built to date, the grounded signal is a human artifact (benchmark suites, unit tests, compilers, proof checkers, game rules). Two boundary cases show why the categorical form of the claim is weaker than the quantitative one, and both are examined in Section~\ref{sec:grounding_discussion}: a conjecture--prove loop running against a fixed proof-checking kernel, and an embodied agent sampling a physical environment with no human in the loop.
\end{remark}

\begin{definition}[Autonomous Self-Improvement]
\label{def:selfimprove}
An autonomous self-improvement process is a sequence of models $\{\bm{\theta}_k\}_{k=0}^{\infty}$ where:
\begin{equation}
\bm{\theta}_{k+1} = \mathcal{A}\!\left(\bm{\theta}_k, \mathcal{D}_k^{\mathrm{self}}\right)
\end{equation}
where $\mathcal{A}$ is a training algorithm and $\mathcal{D}_k^{\mathrm{self}}$ is data generated by the model:
\begin{equation}
\mathcal{D}_k^{\mathrm{self}} = \left\{\bm{x}^{(i)} \sim p_{\bm{\theta}_k}(\cdot)\right\}_{i=1}^{N_k}
\end{equation}
with \textbf{no external grounding signal} (Definition~\ref{def:grounding}), and hence no human influence (Definition~\ref{def:human}), at any step $k \geq 1$. We call this the \emph{sealed loop}.

Two regime parameters matter for the results below and are fixed here. \emph{Replacement}: $\mathcal{D}_k^{\mathrm{self}}$ replaces the previous training corpus rather than being appended to it. \emph{Finite sampling}: $N_k < \infty$, with no smoothing, filtering, or entropy regularization. Section~\ref{sec:regimes} treats the accumulation regime, in which the replacement assumption is dropped, separately---the results differ, and the difference is not cosmetic.
\end{definition}

Earlier versions of this manuscript defined AGI by quantifying over \emph{all} computable functions. That definition is unsatisfiable, and the following proposition---rather than the assumption it replaces---is the reason the criterion below is relativized to a fixed family.

\begin{proposition}[Vacuity of the unrestricted criterion]
\label{prop:vacuity}
Let $\bm{\theta}$ be any model whose inference procedure is computable, and consider the criterion
\begin{equation}
\forall f \text{ computable}, \; \forall \epsilon > 0: \quad
\Pr_{\bm{x} \sim \mathcal{D}_f}\!\left[\argmax_{y} p_{\bm{\theta}}(y \mid \bm{x}) = f(\bm{x})\right] \geq 1 - \epsilon .
\end{equation}
No $\bm{\theta}$ satisfies it. Consequently the set of models meeting the unrestricted criterion is empty, and any claim of the form ``process $P$ cannot converge to it'' is vacuously true of every process, including human-guided ones.
\end{proposition}

\begin{proof}
Fix $\bm{\theta}$. Since inference is computable, the map $\bm{x} \mapsto \argmax_{y} p_{\bm{\theta}}(y \mid \bm{x})$ (with ties broken by a fixed computable rule) is computable. Define
\begin{equation}
g(\bm{x}) := \sigma\!\left(\argmax_{y} p_{\bm{\theta}}(y \mid \bm{x})\right),
\end{equation}
where $\sigma$ is any fixed computable map with no fixed point---appending a symbol, $\sigma(y) = y\,a$, will do. Then $g$ is total computable and $g(\bm{x}) \neq \argmax_y p_{\bm{\theta}}(y \mid \bm{x})$ for every $\bm{x}$. Taking $f = g$, the success probability is $0$, which fails $\geq 1 - \epsilon$ for any $\epsilon < 1$.
\end{proof}

\begin{remark}[What Proposition~\ref{prop:vacuity} does and does not show]
It is a diagonalization against a \emph{fixed} model, so it does not show that intelligence is impossible; it shows that ``matches every computable function'' is the wrong formal target. The same argument defeats the criterion for any finite system, humans included. A target that no process can reach cannot separate autonomous processes from human-guided ones, which is the separation this manuscript is about. This upgrades the v4 nontriviality \emph{assumption} to a proved necessity and is the reason for the relativized definition below.
\end{remark}

\begin{definition}[Artificial General Intelligence --- Operational, Relativized]
\label{def:agi}
Fix, \emph{independently of and prior to} the model under evaluation:
\begin{enumerate}
    \item a countable family $\mathcal{F}$ of computable functions $f : \Sigma^* \to \Sigma^*$, each with a time bound $t_f(n) \leq \mathrm{poly}(n)$ and a description length bound;
    \item for each $f \in \mathcal{F}$ a task distribution $\mathcal{D}_f$ supported on inputs of length at most $n_{\max}$;
    \item a resource budget $(C, b, T)$ for the evaluated model: context length, precision, and inference steps per query.
\end{enumerate}
A model $\bm{\theta}^*$ achieves \emph{AGI relative to $(\mathcal{F}, \{\mathcal{D}_f\}, C, b, T)$} at tolerance $\epsilon > 0$ if
\begin{equation}
\forall f \in \mathcal{F}: \quad
\Pr_{\bm{x} \sim \mathcal{D}_f}\!\left[\argmax_{y} p_{\bm{\theta}^*}(y \mid \bm{x}) = f(\bm{x})\right] \geq 1 - \epsilon .
\end{equation}
The family is fixed before the model, so the diagonal function of Proposition~\ref{prop:vacuity} is not available to the adversary; and the resource bounds make the criterion a statement about achievable accuracy rather than about unbounded computation.
\end{definition}

\begin{proposition}[Nontriviality of the relativized target]
\label{prop:nontrivial}
If $\mathcal{F}$ is finite and every $\mathcal{D}_f$ has finite support, the set of models achieving AGI relative to $(\mathcal{F},\{\mathcal{D}_f\}, C, b, T)$ at tolerance $\epsilon = 0$ is non-empty.
\end{proposition}

\begin{proof}
The set $X = \bigcup_{f \in \mathcal{F}} \mathrm{supp}(\mathcal{D}_f)$ is finite. Define a model that places all its mass on $f(\bm{x})$ for each $\bm{x} \in X$ (a lookup table, realizable as a sufficiently large parameterization at precision $b$, with $|X|$ entries). It matches $f$ with probability~$1$ on every $f \in \mathcal{F}$.
\end{proof}

\begin{remark}[The claim after repair]
\label{rem:agi_repair}
Proposition~\ref{prop:nontrivial} deliberately proves the weakest possible nontriviality: the target is \emph{occupiable}, so the impossibility claim is about \emph{reachability by a sealed autonomous loop}, not about an empty set. It also makes clear what the main theorem does not say: since a lookup table can occupy the target for a finite family, the theorem's force comes from the information accounting of Section~\ref{sec:info}, not from the target being intrinsically unreachable. Extending Proposition~\ref{prop:nontrivial} to infinite $\mathcal{F}$ with resource bounds---where a lookup table is unavailable---remains open and is Priority~1 in the verification plan.
\end{remark}

%% ============================================================
\section{Information-Theoretic Ceiling}
\label{sec:info}
%% ============================================================

\subsection{The Bayesian joint}
\label{sec:joint}

Writing $I(p_{\bm{\theta}}; p_{\mathrm{true}})$ for a \emph{fixed} $p_{\mathrm{true}}$ is not well formed: mutual information between a random object and a constant is zero. Earlier versions of this manuscript left the joint distribution implicit; it is fixed here, and every information quantity below is computed in it.

\begin{definition}[Ground-truth joint]
\label{def:joint}
Let $\mathcal{P}$ be a measurable set of candidate ground-truth distributions over $\Sigma^*$ (equivalently, of candidate answer functions on the task family $\mathcal{F}$ of Definition~\ref{def:agi}), and let $\pi$ be a prior on $\mathcal{P}$. Draw
\begin{equation}
\Theta \sim \pi, \qquad
\mathcal{D}_{\mathrm{train}} \mid \Theta \sim \Theta^{\otimes n}, \qquad
\bm{\theta}_0 = \mathcal{A}_0(\mathcal{D}_{\mathrm{train}}),
\end{equation}
and for $k \geq 0$,
\begin{equation}
\mathcal{D}_k^{\mathrm{self}} \mid \bm{\theta}_k \sim p_{\bm{\theta}_k}^{\otimes N_k}, \qquad
\bm{\theta}_{k+1} = \mathcal{A}(\bm{\theta}_k, \mathcal{D}_k^{\mathrm{self}}),
\end{equation}
where the training algorithm $\mathcal{A}$ may be randomized with randomness independent of $\Theta$. All mutual informations are taken in this joint, and we write $I_k := I(\bm{\theta}_k; \Theta)$.
\end{definition}

Definition~\ref{def:joint} makes the chain
\begin{equation}
\Theta \to \mathcal{D}_{\mathrm{train}} \to \bm{\theta}_0 \to \mathcal{D}_0^{\mathrm{self}} \to \bm{\theta}_1 \to \cdots \to \bm{\theta}_k
\label{eq:chain}
\end{equation}
a genuine Markov chain: conditioned on $\bm{\theta}_k$, the pair $(\mathcal{D}_k^{\mathrm{self}}, \bm{\theta}_{k+1})$ is independent of $\Theta$, because sampling and training in the sealed loop consult nothing else. This conditional independence---not the informal phrase ``self-generated data''---is what the data processing inequality requires, and it is exactly what an external grounding signal (Definition~\ref{def:grounding}) would break.

\begin{theorem}[Information Ceiling]
\label{thm:info_ceiling}
Under Definition~\ref{def:joint} and the sealed-loop condition of Definition~\ref{def:selfimprove}, for every $k \geq 0$:
\begin{equation}
I_k \;\leq\; I_0 \;\leq\; I(\mathcal{D}_{\mathrm{train}}; \Theta) \;\leq\; H(\Theta).
\end{equation}
\end{theorem}

\begin{proof}
Each step of \eqref{eq:chain} is a channel: $\bm{\theta}_{k+1}$ is a function of $(\bm{\theta}_k, \mathcal{D}_k^{\mathrm{self}}, U)$ with $U$ independent of $\Theta$, and $\mathcal{D}_k^{\mathrm{self}}$ is drawn from $p_{\bm{\theta}_k}$. Hence $\Theta \to \bm{\theta}_k \to \bm{\theta}_{k+1}$ is Markov and the data processing inequality gives $I_{k+1} \leq I_k$. Induction gives $I_k \leq I_0$; the same argument applied to $\Theta \to \mathcal{D}_{\mathrm{train}} \to \bm{\theta}_0$ gives $I_0 \leq I(\mathcal{D}_{\mathrm{train}}; \Theta)$, and $I(\cdot;\Theta) \leq H(\Theta)$ always. This is the statement machine-checked in \texttt{Impossibility/InformationCeiling.lean} as \texttt{info\_ceiling}, from the data-processing property carried as a field of the training channel.
\end{proof}

\begin{remark}[No strictness claim]
\label{rem:no_strictness}
Earlier versions asserted \emph{strict} decrease ``whenever $p_{\bm{\theta}_k} \neq p_{\mathrm{true}}$''. That does not follow. Data-processing strictness depends on the channel destroying information---a bijective reparameterization of the parameters destroys none, and $p_{\bm{\theta}_k} \neq p_{\mathrm{true}}$ says nothing about the channel. The conditional core needs only the non-strict inequality, so the strictness claim is withdrawn.
\end{remark}

\begin{remark}[Information vs.\ Computation]
The ceiling is a statement about mutual information in the joint of Definition~\ref{def:joint}. It does not by itself rule out better extraction or use of information already present in the model, and it is agnostic about how much of $I_0$ is computationally accessible (Corollary~\ref{cor:pnp_info}). Information creation and computational extraction are distinct, and only the former is bounded here.
\end{remark}

\subsection{From an information ceiling to an accuracy ceiling}
\label{sec:fano}

Theorem~\ref{thm:info_ceiling} bounds information; Definition~\ref{def:agi} is a statement about \emph{success probability}. The following standard Fano argument supplies the link, which earlier versions of this manuscript left missing: without it, the ceiling constrains a quantity the AGI criterion never mentions.

\begin{lemma}[Information--accuracy link]
\label{lem:fano}
Suppose the prior $\pi$ is supported on $M \geq 2$ candidate ground truths that are pairwise distinguishable on the task family---formally, there is a query $\bm{x}$ and a decoder $W \mapsto f_W(\bm{x})$ separating them---and let $W$ be uniform on those $M$ candidates. Let $\hat{W}$ be any answering rule that is a function of $\bm{\theta}_k$ (and of randomness independent of $\Theta$). Then
\begin{equation}
\Pr[\hat{W} = W] \;\leq\; \frac{I_k + \log 2}{\log M}.
\end{equation}
\end{lemma}

\begin{proof}
$W \to \bm{\theta}_k \to \hat{W}$ is Markov. Fano's inequality gives $H(W \mid \hat{W}) \leq h(P_e) + P_e \log(M-1)$ where $P_e = \Pr[\hat W \neq W]$, and $H(W \mid \hat W) \geq H(W) - I(W; \bm{\theta}_k) = \log M - I_k$ by data processing. Rearranging and bounding $h(P_e) \leq \log 2$ yields $(1 - P_e)\log M \leq I_k + \log 2$.
\end{proof}

\begin{corollary}[Accuracy ceiling under autonomy]
\label{cor:accuracy_ceiling}
Under the hypotheses of Lemma~\ref{lem:fano} and Theorem~\ref{thm:info_ceiling}, every model in the sealed loop satisfies
\begin{equation}
\Pr[\hat{W} = W] \;\leq\; \frac{I_0 + \log 2}{\log M}
\qquad \text{for all } k .
\end{equation}
In particular, if $\log M > \frac{I_0 + \log 2}{1 - \epsilon}$, no iterate meets the AGI criterion of Definition~\ref{def:agi} at tolerance $\epsilon$ on the discriminating queries, and no amount of autonomous iteration changes this: the bound depends on $I_0$ alone.
\end{corollary}

This is the load-bearing chain of the whole program, and it is now machine-checked end to end: sealed loop $\Rightarrow$ information ceiling $\Rightarrow$ accuracy ceiling $\Rightarrow$ failure of the (relativized, nontrivial) AGI criterion. Its cost is explicit: the conclusion is only as strong as the claim that the initial information $I_0$ is small relative to $\log M$, the description complexity of the task family. Making that comparison quantitative---bits of grounding required per unit of capability---is the sharper research question (Section~\ref{sec:grounding_discussion}).

\begin{objection}[Inference-Time Search and Latent Information]
Chain-of-thought reasoning, beam search, and tree search can extract capabilities ``latent'' in the weights without adding new data. Doesn't this violate the information ceiling?
\end{objection}

\begin{response}
No. Inference-time search reorganizes and better utilizes information already present in $\bm{\theta}_k$. This is analogous to lossless compression: you can rearrange bits to make them more useful, but you cannot create bits that were not there. Formally, let $\mathcal{S}$ be any deterministic search procedure applied at inference time. Then:
\begin{equation}
I(\mathcal{S}(p_{\bm{\theta}_k}); p_{\mathrm{true}}) \leq I(p_{\bm{\theta}_k}; p_{\mathrm{true}}) \leq I(\mathcal{D}_{\mathrm{train}}; p_{\mathrm{true}})
\end{equation}
by the data processing inequality (since $\mathcal{S}$ is a function of $p_{\bm{\theta}_k}$). Search can improve \emph{utilization} of existing information up to a ceiling, but that ceiling is fixed by the training data. Better utilization is not the same as new knowledge. The model may solve problems it previously failed on by searching more effectively, but it cannot solve problems whose solutions require information absent from the training distribution. This is the distinction between \emph{extraction} (bounded) and \emph{creation} (what AGI requires for novel domains).

Empirically, Gudibande et al.~\cite{gudibande2023false} demonstrate that fine-tuning open-source LLMs on ChatGPT-generated data transfers surface-level stylistic patterns but fails to close capability gaps on factual and reasoning benchmarks, consistent with the data processing inequality prediction that self-generated data cannot exceed the teacher's information content. Seddik et al.~\cite{seddik2024synthetic} provide a complementary statistical proof that model collapse is unavoidable under purely synthetic training, deriving a critical mixing threshold above which degradation is inevitable regardless of hyperparameters.
\end{response}

%% ============================================================
\section{Distributional Collapse Under Self-Training}
%% ============================================================

A softmax output layer assigns strictly positive probability to every finite sequence. Literal support inclusion $\mathrm{supp}(p_{\bm{\theta}_{k+1}}) \subseteq \mathrm{supp}(p_{\bm{\theta}_k})$ is therefore an equality between full sets for every pair of softmax models: true, trivial, and empty of content. Earlier versions of this manuscript stated the collapse claim that way. It is restated here over the $\epsilon$-effective support, where it has content and a proof.

\begin{definition}[$\epsilon$-effective support]
\label{def:eff_support}
For a distribution $p$ over a countable set and $\epsilon \in (0,1)$, let
\begin{equation}
S_\epsilon(p) := \argmin\left\{ |A| \;:\; A \subseteq \Sigma^*,\; p(A) \geq 1 - \epsilon \right\},
\end{equation}
the smallest set carrying all but $\epsilon$ of the mass, and write $s_\epsilon(p) := |S_\epsilon(p)|$.
\end{definition}

\begin{lemma}[Total-variation transfer]
\label{lem:tv_transfer}
If $\mathrm{TV}(p, q) \leq \eta$ then $s_{\epsilon + \eta}(p) \leq s_{\epsilon}(q)$ for every $\epsilon \in (0, 1-\eta)$.
\end{lemma}

\begin{proof}
Let $A = S_\epsilon(q)$, so $q(A) \geq 1-\epsilon$ and $|A| = s_\epsilon(q)$. Then $p(A) \geq q(A) - \mathrm{TV}(p,q) \geq 1 - \epsilon - \eta$, so $A$ is a witness set for level $\epsilon + \eta$ and the minimal such set is no larger.
\end{proof}

\begin{theorem}[Effective-Support and Entropy Contraction Under Self-Distillation]
\label{thm:collapse}
Let $p_{\bm{\theta}_{k+1}}$ be obtained by fitting $N_k$ samples drawn from $p_{\bm{\theta}_k}$, with no external filtering, selection, smoothing, entropy regularization, or fresh external data, and suppose training reaches $\mathrm{TV}(p_{\bm{\theta}_{k+1}}, \hat{p}_k) \leq \eta$ against the empirical distribution $\hat{p}_k$. Then for every $\epsilon \in (0, 1-\eta)$:
\begin{equation}
s_{\epsilon + \eta}(p_{\bm{\theta}_{k+1}}) \;\leq\; N_k ,
\end{equation}
and, with probability at least $1 - \sqrt{\epsilon}$ over the draw of $\mathcal{D}_k^{\mathrm{self}}$,
\begin{equation}
s_{\sqrt{\epsilon} + \eta}(p_{\bm{\theta}_{k+1}}) \;\leq\; s_\epsilon(p_{\bm{\theta}_k}) .
\end{equation}
Under the same assumptions,
\begin{equation}
H(p_{\bm{\theta}_{k+1}}) \leq H(p_{\bm{\theta}_k}) .
\end{equation}
Consequently, along a sealed loop with per-round budget $N_k \leq N$, the effective support never exceeds $N$ after the first round, regardless of model size.
\end{theorem}

\begin{proof}
\textbf{Sample-budget bound.} $\hat{p}_k$ is supported on at most $N_k$ distinct sequences, so $s_{\epsilon'}(\hat p_k) \leq N_k$ for every $\epsilon' > 0$; Lemma~\ref{lem:tv_transfer} transfers this to $p_{\bm{\theta}_{k+1}}$ at level $\epsilon + \eta$.

\textbf{Teacher bound.} Let $A = S_\epsilon(p_{\bm{\theta}_k})$, so $p_{\bm{\theta}_k}(A^c) \leq \epsilon$. Then $\mathbb{E}[\hat p_k(A^c)] \leq \epsilon$, so by Markov's inequality $\Pr[\hat p_k(A^c) > \sqrt{\epsilon}] \leq \sqrt{\epsilon}$. On the complementary event $\hat p_k(A) \geq 1 - \sqrt{\epsilon}$, hence $s_{\sqrt\epsilon}(\hat p_k) \leq |A| = s_\epsilon(p_{\bm{\theta}_k})$, and Lemma~\ref{lem:tv_transfer} transfers this to $p_{\bm{\theta}_{k+1}}$ at level $\sqrt{\epsilon} + \eta$.

\textbf{Tail pruning is what drives it.} Any sequence with $p_{\bm{\theta}_k}(\bm{x}) < 1/N_k$ is absent from $\mathcal{D}_k^{\mathrm{self}}$ with probability $(1 - p_{\bm{\theta}_k}(\bm{x}))^{N_k} \geq (1 - 1/N_k)^{N_k} \to e^{-1}$: at any fixed budget a constant fraction of the tail is missing from the training signal, and nothing in the sealed loop restores it.

\textbf{Entropy.} Fitting $\hat p_k$ minimizes $\KL(\hat p_k \| p_{\bm{\theta}})$, whose minimizer over a sufficiently expressive class is $\hat p_k$ itself, with $H(\hat p_k) \leq \log N_k$. Entropy contraction therefore holds in the same limit in which the effective-support bound holds; unlike the support bound it is not exact at finite $\eta$, and it is stated here as a consequence of the stated regime rather than as an unconditional inequality.
\end{proof}

\begin{corollary}[Loss of Generality]
\label{cor:loss_generality}
Achieving the criterion of Definition~\ref{def:agi} requires assigning top probability to $f(\bm{x})$ for every $f \in \mathcal{F}$ and $\bm{x} \in \mathrm{supp}(\mathcal{D}_f)$, which requires $s_\epsilon(p_{\bm{\theta}^*}) \geq |\{f(\bm{x})\}|$ for small $\epsilon$. Since the sealed loop caps the effective support at the per-round sample budget, a loop whose budget is below that count moves away from the target rather than toward it.
\end{corollary}

\begin{remark}[What the softmax objection cost the argument]
The repair is not free. The contraction now depends on the training tolerance $\eta$ and the sample budget $N_k$, both of which are engineering parameters: a loop that samples enough, with a large enough $\eta$-slack, contracts slowly. The claim ``self-training collapses'' is therefore quantitative, not categorical, and the interesting regime is where $N_k$ is small relative to the effective support the target requires.
\end{remark}

\begin{objection}[Filtering and Reinforcement Learning]
AlphaZero improved through self-play. If a model generates 100 solutions and trains only on the best one, entropy might decrease but correctness increases. Doesn't this break the collapse theorem?
\end{objection}

\begin{response}
This objection confirms rather than refutes our thesis. Consider what ``selecting the best solution'' requires:
\begin{enumerate}
    \item A \emph{verifier} or \emph{reward function} that evaluates solution quality.
    \item This verifier must encode knowledge of what constitutes a ``good'' solution.
    \item This knowledge was designed by humans.
\end{enumerate}
In AlphaZero~\cite{silver2017alphazero}, the verifier is the game rules---a compact, perfect, human-designed evaluation function. In RLHF~\cite{ouyang2022rlhf}, the reward model is trained on human preferences. In best-of-$N$ sampling with a code interpreter, the interpreter was human-engineered.

Every filtering mechanism that breaks distributional collapse does so by introducing an external signal. By Definition~\ref{def:human}, this constitutes human influence. The target theorem explicitly addresses the case where no such external signal exists. If filtering is needed to avoid collapse, then the role and origin of that filtering signal must be accounted for.

Formally, let $\mathcal{R}: \mathcal{V}^* \to \mathbb{R}$ be a reward function and define the filtered distribution:
\begin{equation}
p_{\bm{\theta}_k}^{\mathcal{R}}(\bm{x}) \propto p_{\bm{\theta}_k}(\bm{x}) \cdot \exp(\beta \cdot \mathcal{R}(\bm{x}))
\end{equation}
The information in $p_{\bm{\theta}_k}^{\mathcal{R}}$ is bounded by:
\begin{equation}
I(p_{\bm{\theta}_k}^{\mathcal{R}}; p_{\mathrm{true}}) \leq I(p_{\bm{\theta}_k}; p_{\mathrm{true}}) + I(\mathcal{R}; p_{\mathrm{true}})
\end{equation}
The term $I(\mathcal{R}; p_{\mathrm{true}})$ is strictly positive---the reward function contributes new information about ground truth. But this information originates from human design. Remove the human-designed $\mathcal{R}$ and you are back to unfiltered self-distillation, which collapses by Theorem~\ref{thm:collapse}.

This collapse has been extensively validated empirically for self-consuming or heavily synthetic training regimes. Alemohammad et al.~\cite{alemohammad2023mad} demonstrate across multiple architectures that self-consuming generative loops cause monotonic increases in FID and decreases in both precision and recall---only fresh external data prevents this entropy contraction. Dohmatob et al.~\cite{dohmatob2024tails} show that iterative self-training on Llama2-7B causes linear degradation in test loss per generation, with low-probability distributional tails vanishing first---precisely the tail-pruning mechanism of Step~1 above. In a stronger result, Dohmatob et al.~\cite{dohmatob2024strong} prove that even 1\% synthetic data contamination triggers collapse, and that larger models amplify rather than mitigate the effect, ruling out the hypothesis that scale alone could overcome distributional collapse. These findings complement the concurrent results of Shumailov et al.~\cite{shumailov2024model}.
\end{response}

%% ============================================================
\section{G\"{o}delian Limits on Self-Verification}
%% ============================================================

\begin{theorem}[Impossibility of Complete Autonomous Self-Verification]
\label{thm:godel}
Let $\mathcal{M}_k$ denote the Turing machine equivalent of $\bm{\theta}_k$. The model cannot completely verify that its self-modifications constitute genuine global improvements across all possible inputs.
\end{theorem}

\begin{proof}
Define $\mathrm{IMPROVE}(k) \iff \mathcal{Q}(\bm{\theta}_{k+1}) > \mathcal{Q}(\bm{\theta}_k)$ for any quality metric $\mathcal{Q}$ correlating with progress toward AGI.

\textbf{G\"{o}del's First Incompleteness Theorem:} $\mathcal{M}_k$, being a formal system capable of arithmetic, contains true statements it cannot prove:
\begin{equation}
\exists\; \phi_k: \quad \mathcal{M}_k \nvdash \phi_k \quad \text{and} \quad \mathcal{M}_k \nvdash \neg\phi_k
\end{equation}
Statements about whether $\bm{\theta}_{k+1}$ will fail on inputs outside $\mathrm{supp}(\mathcal{D}_k^{\mathrm{self}})$ may be among the undecidable statements.

\textbf{G\"{o}del's Second Incompleteness Theorem:} The model cannot prove its own consistency:
\begin{equation}
\mathcal{M}_k \nvdash \mathrm{Con}(\mathcal{M}_k)
\end{equation}

\textbf{Rice's Theorem:} Any non-trivial semantic property of the function computed by $\mathcal{M}_{k+1}$ is undecidable by any Turing machine.

We emphasize that these results do not prevent \emph{all} verification---a model can check specific test cases, run benchmarks, and verify particular outputs. What they prevent is \emph{complete} verification: the model cannot guarantee that $\bm{\theta}_{k+1}$ is globally better than $\bm{\theta}_k$ across all possible inputs. Since AGI requires correctness across all computable functions (Definition~\ref{def:agi}), partial verification is insufficient. The model can confirm improvement on tested cases while unknowingly degrading on untested ones---and it cannot determine, from within itself, whether this is happening.
\end{proof}

\begin{corollary}[The Verification Regress]
\label{cor:regress}
If $\mathcal{M}_k$ builds a verifier $\mathcal{V}_k$, then $\mathcal{V}_k$ is itself a formal system subject to the same incompleteness. Verifying $\mathcal{V}_k$ requires $\mathcal{V}_k'$, generating an infinite regress:
\begin{equation}
\mathcal{M}_k \to \mathcal{V}_k \to \mathcal{V}_k' \to \mathcal{V}_k'' \to \cdots
\end{equation}
Each link in the chain inherits the incompleteness of its predecessor and adds its own statistical error.
\end{corollary}

Empirical evidence supports a weaker practical concern: current LLMs are unreliable self-verifiers without external feedback. Huang et al.~\cite{huang2023selfcorrect} demonstrate that intrinsic self-correction (without external feedback) consistently degrades LLM reasoning performance---the model is more likely to flip a correct answer to incorrect than vice versa. Zheng et al.~\cite{zheng2023judge} document systematic self-enhancement bias in LLM-as-judge evaluations, with GPT-4 favoring its own outputs by 10\%. Most directly, Stechly et al.~\cite{stechly2024verification} show that GPT-4's self-critique causes performance to \emph{decrease} on reasoning tasks, and that the content of critiques is irrelevant---indicating the model cannot reliably generate useful verification signals about its own work.

\subsection{Human Symmetry Objection and Empirical Background}

\begin{objection}[The Human Symmetry Argument]
If G\"{o}delian limits prevent machines from self-improving, they should equally prevent humans, since humans are also finite systems. But humans clearly self-improve (from pre-linguistic hominids to mathematicians). Therefore the G\"{o}delian argument proves too much.
\end{objection}

\begin{response}
This objection identifies a real boundary of the formal argument. The formal proof path does not require proving that humans escape G\"{o}delian limits. It requires only the narrower claim that LLM self-modifications cannot be completely certified by the same formal/computational process that generates them. The following human-brain discussion is therefore empirical and philosophical background, not a theorem in the main proof.

\textbf{The LLM is modeled as a formal system.} A finite-precision LLM consists of:
\begin{itemize}
    \item Finite-precision floating-point arithmetic ($\text{float16}$, $\text{bfloat16}$, or $\text{float32}$),
    \item Deterministic matrix multiplications, additions, and nonlinear activations,
    \item Pseudo-random number generation for sampling (algorithmically deterministic given a seed).
\end{itemize}
Every operation is computable under the finite-precision implementation model. The entire system is realizable as a Turing machine for purposes of the verification-limit argument.

\textbf{The human-brain asymmetry is not established here as a formal theorem.} The brain is a quantum-biological system:
\begin{itemize}
    \item \textbf{Synaptic transmission} involves probabilistic vesicle release governed by quantum tunneling of Ca$^{2+}$ ions through voltage-gated channels.
    \item \textbf{Ion channel dynamics} are inherently quantum mechanical---channel state transitions are stochastic processes with irreducible quantum uncertainty.
    \item \textbf{Molecular processes} underlying neural computation (protein folding, neurotransmitter binding, enzymatic activity) operate in the quantum regime.
\end{itemize}
If any of these quantum processes contribute to cognitive state transitions, then the brain may have access to genuine ontological randomness---not pseudo-randomness generated by an algorithm, but irreducible indeterminacy as characterized by quantum mechanics. Randomness alone, however, does not solve undecidability or establish reliable access to mathematical truth.

\textbf{Conditional physical consequence:} Let $\hat{\rho}_t$ be the quantum density operator describing the brain's state. If the decision process depends on quantum measurement outcomes in a way that cannot be simulated by any effective procedure, then the brain would not be modeled by the same class of deterministic finite-precision computations as deployed LLM inference. This is a physical hypothesis, not a premise discharged by the formal proof.

\textbf{The asymmetry:}
\begin{equation}
\underbrace{\text{LLM}}_{\text{formal system}} \quad \xrightarrow{\text{G\"{o}del applies}} \quad \text{incomplete, cannot fully self-verify}
\end{equation}
\begin{equation}
\underbrace{\text{Human brain}}_{\text{quantum-biological}} \quad \xrightarrow{\text{separate empirical question}} \quad \text{not used as a formal premise}
\end{equation}

This does not require accepting Penrose's specific microtubule hypothesis~\cite{penrose1994}, though the updated Orchestrated Objective Reduction (Orch OR) review~\cite{hameroff2014orchor} provides a detailed theoretical framework. It requires only the weaker claim that quantum effects at the molecular level influence neural computation---and recent experimental evidence now supports this across three independent lines of inquiry.

\textbf{Microtubule quantum effects:} Babcock et al.~\cite{babcock2024superradiance} report the first observation of collective quantum optical superradiance in microtubule tryptophan networks at room temperature, involving $>10^5$ coupled chromophores---an inherently quantum phenomenon. Kalra et al.~\cite{kalra2023exciton} measure electronic energy migration in microtubules at five times the distance predicted by classical F\"{o}rster theory, with clinically used anesthetics significantly dampening this quantum transport. Wiest et al.~\cite{wiest2024epothilone} demonstrate bidirectional causality: microtubule-stabilizing drugs delay anesthetic-induced unconsciousness in rats, establishing that the quantum state of microtubules is functionally linked to consciousness.

\textbf{Nuclear spin effects:} Li et al.~\cite{li2018xenon} demonstrate that xenon isotopes with nuclear spin are approximately 30\% less potent anesthetics than spin-0 isotopes in mice. Since xenon is chemically inert and the only difference between isotopes is nuclear spin, this constitutes direct evidence that a quantum property modulates consciousness.

\textbf{Posner molecules:} Straub et al.~\cite{straub2025posner} provide the first experimental support for Fisher's Posner molecule hypothesis, demonstrating lithium isotope-dependent effects on calcium phosphate formation that are ``entirely unexpected from classical chemistry'' but match quantum dynamical selection predictions (published in PNAS).

These results suggest that quantum effects may be relevant to neural function, but they do not by themselves prove non-computable cognition or an escape from formal limits.

Therefore, the fact that humans can self-improve does not by itself refute the formal self-verification limitation for LLMs. The stronger claim that humans avoid all analogous limits is outside the formal proof path and should be evaluated separately.
\end{response}

%% ============================================================
\section{The SGD Fixed-Point Trap}
%% ============================================================

\begin{theorem}[Self-Training Converges to a Trivial Fixed Point]
\label{thm:sgd_fixed}
Under self-training, SGD converges to a fixed point where the model reproduces its own output distribution.
\end{theorem}

\begin{proof}
The self-training loss is:
\begin{equation}
\mathcal{L}_k(\bm{\theta}) = -\mathbb{E}_{\bm{x} \sim p_{\bm{\theta}_k}}\!\left[\frac{1}{T}\sum_{t=1}^{T} \log p_{\bm{\theta}}(x_t \mid x_{<t})\right] = H(p_{\bm{\theta}_k}) + \KL(p_{\bm{\theta}_k} \| p_{\bm{\theta}})
\end{equation}
This is minimized when $p_{\bm{\theta}} = p_{\bm{\theta}_k}$ (where $\KL = 0$). Therefore $\bm{\theta}_k$ is already a global minimum of $\mathcal{L}_k$. SGD converges back to (approximately) $\bm{\theta}_k$.

The Fisher information matrix captures the effective dimensionality of optimization:
\begin{equation}
\bm{F}_k = \mathbb{E}_{\bm{x} \sim p_{\bm{\theta}_k}}\!\left[\nabla_{\bm{\theta}} \log p_{\bm{\theta}}(\bm{x})\, \nabla_{\bm{\theta}} \log p_{\bm{\theta}}(\bm{x})^\top\right]\bigg|_{\bm{\theta} = \bm{\theta}_k}
\end{equation}
As $p_{\bm{\theta}_k}$ concentrates (Theorem~\ref{thm:collapse}), $\rank(\bm{F}_k)$ decreases. The model's gradients become confined to a shrinking subspace. Optimization can only rearrange existing capabilities within this subspace, not generate novel ones.
\end{proof}

\begin{remark}
This formalizes the intuition: a model training on its own output is trying to become better at being itself. The loss function rewards reproducing what it already does. This is not self-improvement; it is self-reinforcement. Furlanello et al.~\cite{furlanello2018born} demonstrate this empirically: Born Again Networks---iterative self-distillation with identical architectures---converge to a performance ceiling after initial generations. Mobahi et al.~\cite{mobahi2020selfdistillation} prove that self-distillation progressively reduces the number of basis functions in the learned model, converging to a fixed point of minimal representational complexity---providing the formal mechanism underlying this convergence.
\end{remark}

%% ============================================================
\section{The Pattern Completion Barrier}
%% ============================================================

\begin{theorem}[Separation Between Pattern Completion and General Reasoning]
\label{thm:separation}
There exist computable functions $f \in \mathcal{F}_{\mathrm{AGI}}$ that are not approximable by statistical pattern completion on any finite training set.
\end{theorem}

\begin{proof}
We exhibit three classes of functions that are computable but unreachable by pattern matching:

\textbf{(a) G\"{o}del-undecidable theorems:} For any consistent formal system $\mathcal{F}$ (including the implicit system represented by all mathematical text in $\mathcal{D}_{\mathrm{train}}$), there exist true mathematical statements not provable within $\mathcal{F}$. Their proofs cannot be found by pattern-matching on existing proofs.

\textbf{(b) Kolmogorov-complex algorithms~\cite{kolmogorov1965}:} Let $f$ be a function with Kolmogorov complexity $K(f) > K_{\max}(\mathcal{D}_{\mathrm{train}})$. Then $f$ cannot be reconstructed from patterns in $\mathcal{D}_{\mathrm{train}}$.

\textbf{(c) Empirical facts:} Knowledge about the physical world not recorded in $\mathcal{D}_{\mathrm{train}}$ and not logically derivable from it requires observation, which a closed-loop text system cannot perform.
\end{proof}

\begin{objection}[Compositionality]
Even if a program has complexity greater than any single training example, it may be derivable by combining patterns. Humans learn math from examples but prove new theorems.
\end{objection}

\begin{response}
Compositionality allows \emph{recombination} of existing patterns, which is a form of search over the space of combinations. This is addressed by the information ceiling (Theorem~\ref{thm:info_ceiling}): the total information available for composition is bounded by the training data. Composition can extract latent information more efficiently, but it cannot exceed the information-theoretic ceiling. Dziri et al.~\cite{dziri2023faith} demonstrate this directly: transformers solve compositional tasks via linearized subgraph matching rather than systematic reasoning, with accuracy collapsing to near-zero at increasing compositional depth and 82.3\% of correct final answers containing intermediate errors. Berglund et al.~\cite{berglund2023reversal} discover the Reversal Curse---LLMs trained on ``A is B'' achieve near-zero accuracy on ``B is A''---a trivial logical deduction that lies outside what autoregressive pattern completion can extract.

The human analogy remains outside the formal proof path. A human mathematician proving a new theorem may use cognitive processes and external practices that are not captured by finite-context pattern completion, but this manuscript does not require proving that those processes are non-computable.

A transformer is \emph{theoretically} Turing-complete only with unbounded computation time and unbounded memory. In practice, it has a fixed context window $C$ (typically $\leq 2^{17}$ tokens) and fixed precision, making it a finite automaton, strictly weaker than a Turing machine.
\end{response}

%% ============================================================
\section{The Computational Complexity Barrier (P $\neq$ NP)}
%% ============================================================

\begin{theorem}[Self-Improvement Requires Solving NP-Hard Problems]
\label{thm:pnp}
Each step of autonomous self-improvement requires the model to solve an optimization problem over its parameter space that is NP-hard in general. If $\mathrm{P} \neq \mathrm{NP}$, no polynomial-time algorithm (including the model itself) can reliably find optimal or near-optimal solutions.
\end{theorem}

\begin{proof}
At each self-improvement step $k$, the model must find:
\begin{equation}
\bm{\theta}_{k+1}^* = \argmin_{\bm{\theta} \in \mathbb{R}^p} \mathcal{L}_{\mathrm{true}}(\bm{\theta})
\end{equation}
where $\mathcal{L}_{\mathrm{true}}$ is the loss relative to ground truth. This is a non-convex optimization problem over $\mathbb{R}^p$ with $p \sim 10^{9}$ to $10^{12}$ dimensions. Non-convex optimization over continuous spaces with discrete structural constraints is NP-hard in general, as is the closely related problem of training neural networks to global optimality.

The key structural feature is the asymmetry between verification and search:
\begin{itemize}
    \item \textbf{Verification} (checking if a given $\bm{\theta}_{k+1}$ improves on $\bm{\theta}_k$): compute $\mathcal{L}(\bm{\theta}_{k+1})$ and compare. This requires $O(p \cdot C \cdot L)$ operations---polynomial in model size.
    \item \textbf{Search} (finding a $\bm{\theta}_{k+1}$ that improves on $\bm{\theta}_k$): explore a space of size $O(2^{p \cdot b})$ (all possible parameter configurations at $b$-bit precision). In the worst case, this is exponential.
\end{itemize}
This is precisely the $\mathrm{P}$ vs.\ $\mathrm{NP}$ structure: verification is polynomial, but search is (under $\mathrm{P} \neq \mathrm{NP}$) superpolynomial.

Stochastic gradient descent provides a heuristic that avoids exhaustive search by following local gradients. However, SGD:
\begin{enumerate}
    \item Converges to local minima, not global minima,
    \item Cannot escape basins of attraction without external perturbation,
    \item On self-generated data, converges to the trivial fixed point $\bm{\theta}_k$ (Theorem~\ref{thm:sgd_fixed}).
\end{enumerate}
Finding the global optimum---or even a substantially better local optimum---requires solving the NP-hard search problem that SGD explicitly avoids.
\end{proof}

\begin{remark}[Empirical Observations on LLMs and NP-Hard Problems]
\label{rem:pnp_empirical}
An LLM trained on the entirety of available human text has been exposed to billions of instances of solved combinatorial problems---traveling salesman solutions, scheduling solutions, graph coloring solutions, satisfiability instances, cryptographic challenges. If polynomial-time statistical shortcuts existed for NP-hard problems---that is, if there were learnable patterns in the structure of NP solutions---then a sufficiently large pattern-matching system trained on sufficiently many examples would discover them.

Current LLMs are very large pattern-matching systems trained on very large corpora. Despite this:
\begin{enumerate}
    \item LLMs fail to solve large instances of the traveling salesman problem.
    \item LLMs cannot factor large semiprimes.
    \item LLMs cannot find optimal schedules for non-trivial constraint satisfaction problems.
    \item LLMs cannot reliably solve novel SAT instances beyond small sizes.
\end{enumerate}
These failures are evidence about current architectures, training methods, and inference procedures. They are not a proof of $\mathrm{P} \neq \mathrm{NP}$ and should not be used as a theorem-level premise. The formal complexity barrier remains conditional on the relevant optimization problem and complexity assumptions.
\end{remark}

\begin{corollary}[P $\neq$ NP Strengthens the Information Ceiling]
\label{cor:pnp_info}
Even information that is technically ``latent'' in the training data (the information ceiling of Theorem~\ref{thm:info_ceiling}) may be computationally inaccessible. Under $\mathrm{P} \neq \mathrm{NP}$, extracting certain latent information requires superpolynomial computation:
\begin{equation}
\exists\; f \in \mathrm{NP}: \quad f \text{ is derivable from } \mathcal{D}_{\mathrm{train}} \text{ but requires } 2^{\Omega(n)} \text{ steps to extract}
\end{equation}
Therefore, the effective information available to the model is strictly less than the theoretical information ceiling:
\begin{equation}
I_{\mathrm{accessible}}(\bm{\theta}_k; p_{\mathrm{true}}) < I(\bm{\theta}_k; p_{\mathrm{true}}) \leq I(\mathcal{D}_{\mathrm{train}}; p_{\mathrm{true}})
\end{equation}
The information ceiling was already a hard limit. Computational complexity makes the \emph{effective} ceiling even lower.
\end{corollary}

\begin{remark}[The Self-Improvement Loop as Iterated Hard Search]
Under the stated complexity assumptions, each iteration of autonomous self-improvement requires solving a hard optimization problem (finding better parameters). The model's tool for this search is statistical pattern matching. Current failures on NP-hard benchmarks are empirical corroboration of limited search ability, not proof of $\mathrm{P} \neq \mathrm{NP}$. Fan et al.~\cite{fan2023nphardeval} benchmark 12 LLMs across complexity classes, finding weighted accuracy collapses from 0.25 (NP-complete) to 0.02 (NP-hard). If $\mathrm{P} \neq \mathrm{NP}$ and the parameter-search problem instantiates an NP-hard family, additional parameters, data, or compute do not by themselves convert polynomial verification into guaranteed superpolynomial search.
\end{remark}

%% ============================================================
\section{Resource Bounds on In-Context Self-Modification}
\label{sec:context}
%% ============================================================

\begin{remark}[Scope of this section]
\label{rem:context_scope}
The results below bound one particular self-improvement strategy: the model representing its own parameters \emph{in its context window} and reasoning about them as tokens. Definition~\ref{def:selfimprove} does not require that strategy. The loop it defines updates parameters through the training algorithm $\mathcal{A}$, and gradient descent touches every parameter without any parameter ever entering the context. Treating this section as a barrier to autonomous self-improvement in general would be attacking a scenario the formal loop never needs; it is retained because in-context self-modification is a strategy actually proposed in the literature, and because the bounds are quantitative and checkable. The sharper form of the underlying worry is introspective rather than architectural: a system that cannot simulate its own operation cannot target modifications at it~\cite{zhang2026introspection}.
\end{remark}

\begin{theorem}[In-Context Self-Correction Requires Unbounded Resources]
\label{thm:context}
Let $\bm{\theta}_k \in \mathbb{R}^p$ be the model parameters at step $k$, with $p$ on the order of $10^{9}$ to $10^{12}$. Let $C$ be the context window size (maximum number of tokens processable in a single forward pass). For the model to verify or correct its own parameters, it must represent and reason about the state space $\mathbb{R}^p$. The following bound holds:
\begin{equation}
\frac{p \cdot b}{C \cdot \log_2 |\mathcal{V}|} \gg 1
\end{equation}
where $b$ is the number of bits per parameter (typically 16 or 32). That is, the information required to describe the model's own parameter space vastly exceeds the model's processing capacity in any single reasoning step.
\end{theorem}

\begin{proof}
A model with $p = 10^{11}$ parameters at 16-bit precision contains $1.6 \times 10^{12}$ bits of state information. With a context window of $C = 128{,}000$ tokens and vocabulary size $|\mathcal{V}| \approx 10^5$, each context window carries approximately $128{,}000 \times 17 \approx 2.2 \times 10^6$ bits of information.

The ratio:
\begin{equation}
\frac{\text{model state}}{\text{context capacity}} \approx \frac{1.6 \times 10^{12}}{2.2 \times 10^{6}} \approx 7.3 \times 10^5
\end{equation}
The model would need approximately $730{,}000$ sequential reasoning steps to even \emph{read} its own parameters, let alone reason about them.

For self-correction, the model must:
\begin{enumerate}
    \item Identify which parameters contribute to errors (requires gradient computation across the full parameter space),
    \item Determine correct parameter values (requires optimization over $\mathbb{R}^p$),
    \item Verify that corrections do not introduce new errors (requires evaluation across all possible inputs).
\end{enumerate}
Each step requires processing information proportional to $p$ or larger. Since the model can process at most $C$ tokens per step, the number of required steps grows as $\Omega(p / C)$ at minimum. For verification across all inputs (step 3), the required computation is unbounded.

Therefore, even if the model could in principle self-correct (which Theorem~\ref{thm:godel} shows it cannot do completely), the computational resources required grow without bound relative to its finite processing capacity:
\begin{equation}
\lim_{p \to \infty} \frac{\text{Resources required for self-correction}}{\text{Resources available per step}} = \infty
\end{equation}
\end{proof}

\begin{theorem}[The Overhead Paradox]
\label{thm:overhead}
Even in the theoretical best case where the context window $C$ is extended to equal the full parameter representation ($C \cdot \log_2 |\mathcal{V}| = p \cdot b$), self-correction remains impossible due to computational overhead that creates an irreducible deficit.
\end{theorem}

\begin{proof}
Suppose we grant the model a context window large enough to hold its entire parameter state. The context must then be partitioned among competing demands:
\begin{equation}
C = C_{\mathrm{state}} + C_{\mathrm{reasoning}} + C_{\mathrm{temp}} + C_{\mathrm{output}}
\end{equation}
where:
\begin{itemize}
    \item $C_{\mathrm{state}}$: tokens representing the parameters to be examined,
    \item $C_{\mathrm{reasoning}}$: tokens for the logical operations analyzing those parameters (comparisons, gradient estimation, error diagnosis),
    \item $C_{\mathrm{temp}}$: tokens for intermediate results and working storage,
    \item $C_{\mathrm{output}}$: tokens for the corrected parameter specifications.
\end{itemize}
Since $C_{\mathrm{state}} = C$ (by assumption, the model's full state fills the context), we require:
\begin{equation}
C_{\mathrm{reasoning}} + C_{\mathrm{temp}} + C_{\mathrm{output}} > 0
\end{equation}
for any non-trivial reasoning to occur. But $C_{\mathrm{state}} + C_{\mathrm{reasoning}} + C_{\mathrm{temp}} + C_{\mathrm{output}} > C$, which is a contradiction.

In practice, meaningful self-analysis requires substantial reasoning overhead. At minimum, $C_{\mathrm{reasoning}} \geq \alpha \cdot C_{\mathrm{state}}$ for some constant $\alpha > 0$ reflecting the complexity of error diagnosis. Therefore:
\begin{equation}
C_{\mathrm{required}} \geq (1 + \alpha) \cdot C_{\mathrm{state}} + C_{\mathrm{temp}} + C_{\mathrm{output}} > C
\end{equation}
The deficit is irreducible: \emph{any} system reasoning about a representation of itself requires more resources than the representation alone occupies.
\end{proof}

\begin{theorem}[The Compute Cost of Buying Context]
\label{thm:scaling}
Under standard (quadratic) attention, extending the context window to hold a parameter-sized state makes a single self-analysis pass quadratic in the size of that state.
\end{theorem}

\begin{remark}[A premise withdrawn]
\label{rem:rope}
Earlier versions of this manuscript argued that enlarging $C$ does not help because the parameter count $p = g(C)$ grows monotonically with the context window. That premise is empirically false. RoPE-style positional interpolation and its successors extend fixed-size models from 8K to 128K+ tokens with \emph{no} added positional parameters, and the embedding matrix does not scale with $C$ at all in such schemes. What does scale with $C$ is \emph{compute}, not parameter count. The argument below is restated in that form; the claim that a larger context automatically implies a proportionally larger model is withdrawn.
\end{remark}

\begin{proof}
The \emph{computational cost} of processing context $C$ in the standard transformer is:
\begin{equation}
\mathrm{FLOPS}_{\mathrm{attention}} = O(C^2 \cdot d)
\end{equation}
due to the quadratic scaling of self-attention. Doubling $C$ quadruples the compute required per forward pass. The time to perform one self-reasoning step therefore grows as:
\begin{equation}
T_{\mathrm{step}}(C) = O(C^2)
\end{equation}
For a model attempting to reason about its full parameter space ($C \propto p$), the time per self-correction step is:
\begin{equation}
T_{\mathrm{step}} = O(p^2)
\end{equation}
For $p = 10^{11}$, this yields $T_{\mathrm{step}} \propto 10^{22}$ FLOPS per reasoning step. At current hardware speeds ($\sim10^{18}$ FLOPS/s for the largest clusters), a single self-reasoning pass would take approximately $10^4$ seconds $\approx 3$ hours---for one step of what would require millions of sequential steps to cover the parameter space. The bound is architecture-specific: sub-quadratic attention variants weaken the exponent, and the argument should be read as a statement about the standard architecture rather than about all sequence models.
\end{proof}

\begin{theorem}[Training Time Divergence]
\label{thm:training_time}
The time required for the model to retrain itself after identifying corrections grows superlinearly with model size, creating a secondary resource explosion.
\end{theorem}

\begin{proof}
Suppose the model, at self-improvement step $k$, has identified a set of corrections to apply to its parameters. Implementing these corrections requires a training process (gradient-based optimization). A single training step requires:
\begin{equation}
T_{\mathrm{train\_step}} = O(p \cdot C \cdot L)
\end{equation}
FLOPS, where $L$ is the number of transformer layers. Convergence to a new parameter configuration typically requires $S$ training steps, where $S$ depends on the magnitude of the correction and the loss landscape curvature. The total training time is:
\begin{equation}
T_{\mathrm{train\_total}} = S \cdot O(p \cdot C \cdot L)
\end{equation}

For the self-improvement loop, this training must occur at \emph{every} iteration $k$. The total compute for $K$ iterations of self-improvement is:
\begin{equation}
T_{\mathrm{total}} = \sum_{k=0}^{K-1} \left[\underbrace{T_{\mathrm{diagnose}}(k)}_{\text{identify errors}} + \underbrace{T_{\mathrm{train}}(k)}_{\text{apply corrections}} + \underbrace{T_{\mathrm{verify}}(k)}_{\text{check results}}\right]
\end{equation}

By Theorem~\ref{thm:sgd_fixed}, the achievable improvement $\Delta_k \to 0$, meaning each successive iteration produces diminishing returns. But the \emph{cost} per iteration does not diminish---it remains $\Omega(p \cdot C \cdot L)$ per training cycle. The ratio of cost to benefit therefore diverges:
\begin{equation}
\frac{T_{\mathrm{train}}(k)}{\Delta_k} \to \infty \quad \text{as } k \to \infty
\end{equation}

Furthermore, as errors accumulate (Theorem~\ref{thm:error_compound}), the corrections required become more complex (affecting larger portions of the parameter space), increasing $S_k$. Simultaneously, the model's ability to correctly diagnose errors degrades (since it is using a progressively more corrupted model to identify corruption). This creates a feedback loop:
\begin{enumerate}
    \item More errors $\implies$ more parameters need correction $\implies$ longer training,
    \item Longer training on self-generated data $\implies$ more distributional collapse (Theorem~\ref{thm:collapse}),
    \item More collapse $\implies$ worse error diagnosis $\implies$ more errors introduced.
\end{enumerate}

The self-improvement process is therefore not merely slow but \emph{anti-convergent}: the resources required per unit of improvement grow without bound, while the quality of each improvement attempt simultaneously degrades. In the limit:
\begin{equation}
\lim_{k \to \infty} \frac{\text{Compute cost of iteration } k}{\text{Genuine improvement from iteration } k} = \infty
\end{equation}
Under these assumptions, even with unbounded time and compute, the process does not converge---it diverges, because the errors introduced at each step eventually exceed the corrections achieved.
\end{proof}

\begin{remark}[The Infinite Loop]
Theorems~\ref{thm:context}--\ref{thm:training_time} together describe an inescapable recursive trap:
\begin{enumerate}
    \item The model cannot fit a representation of itself into its working memory (\ref{thm:context}).
    \item Even if it could, the overhead of reasoning about that representation exceeds the available space (\ref{thm:overhead}).
    \item Expanding the context does not come free: under standard attention, the compute per self-analysis pass grows quadratically in the state to be analyzed (\ref{thm:scaling}).
    \item The time to retrain after each correction cycle diverges relative to the improvement gained (\ref{thm:training_time}).
\end{enumerate}
Under the resource-scaling assumptions above, each attempt to provide the system with enough working memory for in-context self-correction raises the compute cost of using it. This is a cost curve, not a proof of impossibility: it says in-context self-modification at frontier parameter counts is economically hopeless with quadratic attention, not that no autonomous loop can update parameters---which, again, is what Remark~\ref{rem:context_scope} sets aside. Combined with the conditional error model of Theorem~\ref{thm:error_compound}, this creates an anti-convergent loop rather than a guaranteed path to self-correction.

Empirically, Liu et al.~\cite{liu2023lost} demonstrate that LLMs exhibit a U-shaped performance curve over context position, with significant degradation for information in the middle---if a model cannot reliably access information throughout its own context, the comprehensive self-analysis required for self-improvement is infeasible. The scaling trap is quantified by Kaplan et al.~\cite{kaplan2020scaling}, who show language model loss follows power-law scaling with exponents of $\sim$0.07 for parameters, guaranteeing diminishing returns where each $10\times$ compute increase yields only $\sim$12\% loss reduction. Hoffmann et al.~\cite{hoffmann2022chinchilla} confirm this with compute-optimal training: Chinchilla (70B) outperforms Gopher (280B) by only 7\% on MMLU despite vastly more efficient resource allocation.
\end{remark}

%% ============================================================
\section{Error Compounding in Recursive Self-Modification}
%% ============================================================

\begin{theorem}[Non-Improvement Without an Independent Corrector]
\label{thm:error_compound}
Let $\epsilon_k$ measure the error of $\bm{\theta}_k$ relative to ground truth, and assume the recurrence
\begin{equation}
\epsilon_{k+1} = \epsilon_k - \Delta_k + \delta_k + \sum_{j=0}^{k} \alpha_{k,j}\, \delta_j ,
\label{eq:err_rec}
\end{equation}
where $\Delta_k \geq 0$ is the intended improvement, $\delta_k \geq 0$ is undetected degradation, and $\alpha_{k,j} \geq 0$ captures error propagation through self-generated training data.
\begin{enumerate}
    \item \textbf{(Universal form.)} If $\Delta_k \leq \delta_k$ for all $k$---the corrector and the corrupted system being the same stochastic process---then $\epsilon_{k+1} \geq \epsilon_k$ for all $k$: the error never improves.
    \item \textbf{(Divergence, replacement regime.)} If in addition $\delta_k \geq \delta > 0$ for all $k$ and $\Delta_k \to 0$, then $\epsilon_k \to \infty$; concretely, once $\Delta_k \leq \delta/2$ the error grows at least linearly, $\epsilon_{K+n} \geq \epsilon_K + n\delta/2$.
    \item \textbf{(Stagnation at the fixed point.)} If $\Delta_k = \delta_k = 0$---the exact fixed point of Theorem~\ref{thm:sgd_fixed}---then $\epsilon_{k+1} = \epsilon_k$ for all $k$.
\end{enumerate}
\end{theorem}

\begin{proof}
All three parts follow from \eqref{eq:err_rec} and $\alpha_{k,j}\delta_j \geq 0$. (1) $\epsilon_{k+1} - \epsilon_k = (\delta_k - \Delta_k) + \sum_j \alpha_{k,j}\delta_j \geq 0$. (2) For $k \geq K$, $\epsilon_{k+1} - \epsilon_k \geq \delta - \delta/2 = \delta/2$; induction gives the linear bound, and the Archimedean property gives divergence. (3) Every term after $\epsilon_k$ vanishes. All three are machine-checked in \texttt{Impossibility/ErrorDivergence.lean}.
\end{proof}

\begin{remark}[Which form the main theorem needs]
\label{rem:which_form}
Only part~(1). Convergence to the AGI criterion requires the error to \emph{fall}; a process whose error never falls does not reach it. Divergence is a stronger conclusion that requires the stronger hypotheses of part~(2), and the manuscript no longer asserts it universally.
\end{remark}

\begin{remark}[Reconciling Theorem~\ref{thm:sgd_fixed} and Theorem~\ref{thm:error_compound}]
\label{rem:reconcile}
Earlier versions assumed $\delta_k > 0$ as part of the error model while simultaneously arguing that self-training converges to a fixed point where the model reproduces itself. Those two claims are in tension: at an exact fixed point nothing is being changed, so nothing is being undetectably degraded, and $\epsilon_{k+1} = \epsilon_k$. Part~(3) records this. The resolution is that a real loop sits near, not at, the fixed point: finite-sample noise keeps $\delta_k > 0$ while $\Delta_k \to 0$, which yields drift rather than either exact stagnation or guaranteed divergence. The defensible universal picture is \emph{stagnation with random-walk drift}, matching the entropy-decay and variance-amplification invariants derived independently by Zenil~\cite{limitsself2025}, with divergence as the special case where the drift has a persistent sign.
\end{remark}

\subsection{The accumulation regime}
\label{sec:regimes}

The divergence half of part~(2) is regime-specific, and the regime matters.

\begin{objection}[Accumulation bounds the error]
Gerstgrasser et al.~\cite{gerstgrasser2024accumulate} prove that when synthetic data \emph{accumulate} alongside the original real data instead of replacing it, the test error has a finite upper bound independent of the number of iterations---model collapse does not occur. Their linear-model analysis is accompanied by experiments across transformers, diffusion models, and VAEs. Doesn't this refute the error-divergence claim?
\end{objection}

\begin{response}
It refutes the universal reading, and the theorem above has been restated accordingly. Three points remain.

First, the accumulation result bounds error; it does not produce improvement. A process whose error is bounded but whose information about $\Theta$ is capped by Theorem~\ref{thm:info_ceiling} stagnates below the ceiling---it does not converge to the AGI criterion. The accumulation regime therefore defeats part~(2), not part~(1) or the main theorem.

Second, Definition~\ref{def:selfimprove} explicitly fixes the replacement regime, and this is a real narrowing of scope rather than a technicality. A model that has memorized portions of $\mathcal{D}_{\mathrm{train}}$ and samples at low temperature partially recreates accumulation from inside the sealed loop while formally satisfying ``data sampled from $p_{\bm{\theta}_k}$''. The boundary between the regimes is therefore softer than the definitions suggest, and marking it precisely---in terms of how much of $\mathcal{D}_{\mathrm{train}}$ survives in $\mathcal{D}_k^{\mathrm{self}}$---is open work.

Third, the accumulation result buys boundedness at the price of unbounded storage of the original human corpus: the real data must be retained forever. That is a dependence on the externally grounded corpus, made permanent rather than removed, which is consistent with the framing of Remark~\ref{rem:grounding} but is a much weaker claim than ``autonomous loops diverge''.
\end{response}

\begin{remark}[The Statistical Error Floor]
The fundamental issue is that the system is stochastic. Every output has an associated error probability. When the system uses its own outputs to train itself, it trains on data that contains errors. The next generation inherits those errors and adds its own. The error is cumulative and the correction mechanism (the model itself) has the same statistical error rate as the system it is trying to correct. This is not a bug to be fixed; it is a mathematical property of using statistics to correct statistics.

An external corrector (human or human-designed verifier) breaks this cycle because its error properties are independent of the model's. But an AI-generated corrector has correlated errors---it was built by the same flawed process---and therefore cannot provide the independent error correction needed. Gao et al.~\cite{gao2022overoptimization} establish scaling laws for reward model overoptimization: true quality initially improves then degrades as optimization against any proxy intensifies, demonstrating that even human-designed reward models degrade under iterative optimization---the problem would be far worse with self-generated evaluation. A concurrent formal analysis~\cite{limitsself2025} independently models recursive self-training as a discrete-time dynamical system, proving entropy decay and variance amplification invariants that hold regardless of model architecture.
\end{remark}

%% ============================================================
\section{Main Theorem}
%% ============================================================

\begin{theorem}[Conditional Impossibility of Autonomous LLM Self-Improvement to AGI]
\label{thm:main}
Let $\{p_{\bm{\theta}_k}\}_{k=0}^{\infty}$ be the sequence of models produced by autonomous self-improvement (Definition~\ref{def:selfimprove}) in the joint of Definition~\ref{def:joint}. Then $\{p_{\bm{\theta}_k}\}$ cannot converge to AGI relative to a fixed family (Definition~\ref{def:agi}, nontrivial by Proposition~\ref{prop:nontrivial}) unless the initial model already satisfies the AGI criterion.

The proof needs only barriers~1 and~2 below: the information ceiling with its accuracy link (Corollary~\ref{cor:accuracy_ceiling}), and the sealed-loop condition that makes the chain Markov. This is the version machine-checked in \texttt{Impossibility/MainTheorem.lean}: given $I_0 < H(\Theta)$, convergence would force $I_k \geq H(\Theta) - \epsilon$ for every $\epsilon$, contradicting $I_k \leq I_0$ at $\epsilon = (H(\Theta) - I_0)/2$. The remaining barriers are not load-bearing for the conclusion; they describe additional, independent ways the sealed loop fails, and are listed with their current status.

\begin{enumerate}
    \item \textbf{Information ceiling} (Theorem~\ref{thm:info_ceiling}, Corollary~\ref{cor:accuracy_ceiling}) --- \emph{load-bearing; machine-checked.} In the joint of Definition~\ref{def:joint}, self-generated data cannot create mutual information about $\Theta$, and Fano converts that cap into a cap on achievable accuracy. Inference-time search improves utilization of existing information without raising the ceiling.

    \item \textbf{Distributional collapse} (Theorem~\ref{thm:collapse}) --- \emph{machine-checked from the finite-sample assumptions.} Under unsmoothed, unfiltered self-distillation the $\epsilon$-effective support is capped by the per-round sample budget and entropy does not rise. Stated over effective support, not literal support, which is vacuous for softmax models. Mechanisms that prevent collapse inject an external grounding signal by definition.

    \item \textbf{G\"{o}delian verification failure} (Theorem~\ref{thm:godel}) --- \emph{Rice component machine-checked; G\"{o}del~I in computability form machine-checked; G\"{o}del~II carried as a named bridge axiom.} The model cannot completely verify its self-modifications as genuine global improvements. Partial verification (benchmarks, test cases) is possible and is exactly where external grounding enters. The human-brain asymmetry discussion is empirical background, not a premise.

    \item \textbf{SGD fixed-point trap} (Theorem~\ref{thm:sgd_fixed}) --- \emph{population statement machine-checked; Fisher-rank contraction carried as a hypothesis.} The teacher distribution minimizes the self-distillation objective, so the loop optimizes toward reproducing itself.

    \item \textbf{Pattern completion barrier} (Theorem~\ref{thm:separation}) --- \emph{not formalized; depends on Kolmogorov complexity, which has no Lean~4 formalization.} The target claim is that pattern completion on finite training data is weaker than general computation under precise finite-model assumptions.

    \item \textbf{Computational complexity barrier} (Theorem~\ref{thm:pnp}, Corollary~\ref{cor:pnp_info}, Remark~\ref{rem:pnp_empirical}) --- \emph{conditional; not formalized.} If the search problem instantiates an NP-hard family and $\mathrm{P} \neq \mathrm{NP}$, the verification/search gap is permanent. Current LLM failures on NP-hard benchmarks are empirical context, not evidence for $\mathrm{P} \neq \mathrm{NP}$.

    \item \textbf{Resource bounds on in-context self-modification} (Theorems~\ref{thm:context}--\ref{thm:training_time}) --- \emph{quantitative bounds on one strategy, not on the loop of Definition~\ref{def:selfimprove}} (Remark~\ref{rem:context_scope}). Working memory cannot hold the parameter state; reasoning about a state that fills the context needs more room than the state occupies; buying context costs quadratic compute under standard attention; and cost per unit of improvement diverges as improvement vanishes.

    \item \textbf{Non-improvement} (Theorem~\ref{thm:error_compound}) --- \emph{machine-checked, in the weakened form.} When undetected degradation is at least as large as intended improvement, error never falls. Divergence holds in the replacement regime with a positive degradation floor; in the accumulation regime error is provably bounded (Section~\ref{sec:regimes}), so divergence is not claimed universally.
\end{enumerate}

\noindent The conclusion is:
\begin{equation}
\boxed{\forall\; \bm{\theta}_0 \in \mathbb{R}^p \text{ with } I_0 < H(\Theta),\;\; \forall\; k \geq 0: \quad \mathcal{S}^k(\bm{\theta}_0) \not\to \bm{\theta}^*_{\mathrm{AGI}}}
\end{equation}
within the sealed autonomous stochastic autoregressive framework, excluding the degenerate case in which $\bm{\theta}_0$ already carries AGI-level information about $\Theta$. What remains open is not the derivation but the \emph{interpretation of its hypotheses}: whether the data-processing property of a real training pipeline holds as stated, how large $I_0$ is relative to $\log M$ for realistic task families, and---the question the 2025--26 empirical record actually poses---how much capability a small amount of external grounding unlocks (Section~\ref{sec:grounding_discussion}).
\end{theorem}

%% ============================================================
\section{Discussion}
%% ============================================================

\subsection{What the Proof Program Targets}

The proof program targets the following conditional result: if the information ceiling, collapse, verification, fixed-point, complexity, resource, and error-divergence barriers are formalized under the definitions above, then a stochastic language model operating without any form of human influence cannot self-improve to AGI unless it already satisfies the AGI criterion. The intended scope is broad with respect to model size, context length, and compute budget, but those scaling claims remain part of the verification work.

\subsection{What Would Escape These Bounds}

A system escapes these bounds if it incorporates:
\begin{itemize}
    \item External information sources (new data, environment interaction) --- breaks the sealed-loop conditional independence of Definition~\ref{def:joint}; an embodied agent is the clean case, and it is outside the scope of these results rather than refuted by them,
    \item Evaluation signals with $I(S;\Theta \mid \bm{\theta}_k) > 0$ (reward models, verifiers, unit tests, compilers, proof checkers, game rules) --- external grounding, of which human design is the dominant special case,
    \item Data accumulation rather than replacement --- escapes the collapse and divergence results while retaining a permanent dependence on the original grounded corpus (Section~\ref{sec:regimes}),
    \item Non-stochastic components (formal theorem provers, symbolic reasoning engines) --- departs from the stochastic framework; this is also the direction Zenil~\cite{limitsself2025} argues is necessary.
\end{itemize}
Every empirically demonstrated self-improvement system we are aware of takes the second or third route (Section~\ref{sec:empirical}). What the results here rule out is the sealed case: no external grounding, replacement dynamics, text-only. Whether that is the case the singularity hypothesis actually requires is a substantive question, and the honest answer is that the strongest current systems do not attempt it.

\subsection{The Stochastic Parrot, Formalized}

The proof program aims to provide mathematical grounding for the characterization of LLMs as ``stochastic parrots''~\cite{bender2021stochastic}. The model is a conditional probability estimator: $p_{\bm{\theta}}(x_t \mid x_{<t})$. Self-improvement within this framework is self-distillation: the model learning to reproduce its own distribution. Under the information-ceiling assumptions, the output may become more fluent, more consistent, or more confident without becoming more correct about things absent from training or more general than the information-theoretic ceiling permits.

\subsection{Human Influence as Theorem Target}
\label{sec:human_target}

The central theorem target is not that LLMs cannot improve---they clearly can and do. It is that \emph{human influence is a necessary component of open-ended improvement} under the autonomous stochastic-language-model framework. Whether through curated training data, designed reward functions, engineered tools, or architectural innovations, current LLM progress depends on external structure. The formal task is to show that severing this dependency requires escaping at least one of the assumptions behind the barriers above. Until those assumptions are discharged, this should be read as a conditional impossibility claim rather than a completed theorem.

\subsection{From ``Human Influence'' to ``External Grounding''}
\label{sec:grounding_discussion}

The reformulation announced in Remark~\ref{rem:grounding} has consequences that are worth stating plainly, including where it costs the thesis something.

\paragraph{What the formal arguments actually use.} The data processing inequality is indifferent to the provenance of a signal. Theorem~\ref{thm:info_ceiling} needs the sealed-loop conditional independence; Theorem~\ref{thm:collapse} needs the absence of filtering; Corollary~\ref{cor:accuracy_ceiling} needs the information cap. None of them mentions humans. The correct general statement of the thesis is therefore: \emph{open-ended improvement requires external grounding}, with human design as the empirically dominant, and historically universal, source of it.

\paragraph{Boundary case (a): deduction.} Consider a sealed conjecture--prove loop equipped with a fixed proof-checking kernel: the model proposes statements and candidate proofs, the kernel accepts or rejects, and accepted theorems are added to the training corpus. Such a loop can generate unboundedly many new true theorems, none of which appeared in $\mathcal{D}_{\mathrm{train}}$. Two observations, in tension:

\begin{itemize}
    \item Formally, this does not violate Theorem~\ref{thm:info_ceiling}. The kernel's acceptances are a deterministic function of the conjecture and the axioms, so the loop's information about $\Theta$---about which mathematical universe it is in---does not increase. What increases is the \emph{computationally accessible} portion of information already implied by the axioms, which Corollary~\ref{cor:pnp_info} explicitly separates from the information ceiling.
    \item Practically, this is a serious weakening of the categorical claim. A proof kernel is a small, static human artifact. If a few thousand lines of human-written checking code are enough to unlock open-ended capability growth in a domain, then ``human influence is necessary'' is true but nearly vacuous as a safety or forecasting claim---the required human contribution is a one-time constant, not an ongoing dependence.
\end{itemize}

\paragraph{Boundary case (b): embodiment.} An agent that acts in a physical environment and observes the results obtains signal with $I(S;\Theta \mid \bm{\theta}_k) > 0$ and no human in the loop. Definition~\ref{def:human} does not cover it; Definition~\ref{def:grounding} does. This manuscript excludes such systems only via the closed-loop clause of Definition~\ref{def:selfimprove}, which is a scope restriction rather than an argument. We state explicitly: nothing here shows that an embodied autonomous learner is bounded, and the claim ``human involvement is mathematically necessary for AI progress'' is \emph{not} established by these results---what is established concerns sealed text-only self-distillation.

\paragraph{The quantitative question.} The interesting form of the thesis is an exchange rate: how much capability does a bit of grounded signal buy, and how far can a fixed, finite grounding source carry a loop before it saturates? Three recent results turn this from rhetoric into measurable structure. Feng et al.~\cite{feng2024verification} show that even imperfect verifiers prevent collapse and enable scaling with synthesized data. Yi et al.~\cite{yi2026escaping} analyze verifier-guided retraining and prove that the estimate converges to the \emph{verifier's} knowledge center: early gains are real, and they plateau---or reverse---exactly to the extent the verifier is imperfect. Song et al.~\cite{song2024gap} identify the generation--verification gap as the governing quantity for self-improvement and find it scaling monotonically with pre-training compute, i.e., with the human corpus. Together these sketch the theorem this program should be targeting:

\begin{quote}
\emph{Conjecture (bits-of-grounding).} For a sealed loop with a grounding source $S$, achievable capability is bounded by a function of $I(S;\Theta)$ that saturates at the source's own competence, independently of iteration count, model size, and compute.
\end{quote}

That statement is falsifiable, is consistent with everything in the empirical record of Section~\ref{sec:empirical}, and would carry real forecasting content---which the categorical claim does not.

\subsection{A Control Theory Perspective}
\label{sec:control}

Control theory provides an intuitive framing for the central impossibility. A feedback control system works because the controller is \emph{independent} of the plant. A thermostat is not made of the same heat it regulates.

An LLM attempting to regulate itself is simultaneously the plant and the controller. The sensor, actuator, and error estimator are all the same stochastic process with the same systematic biases. When the system drifts, the component detecting the drift has drifted by the same amount in the same direction.

In control theory terms, the controller's noise is coupled to the plant's noise with correlation approaching~$1$. This violates the fundamental assumption that makes feedback control work: independent error correction.

More precisely, classical control theory requires four prerequisites that LLMs violate:

\begin{enumerate}
    \item \textbf{Observable state.} The meaningful state of an LLM is the distribution $p_{\bm{\theta}}$ over all possible inputs---a function over an exponentially large space. This state is not observable; it can only be sampled. No finite number of samples characterizes the distribution.
    \item \textbf{Plant model.} The relationship between parameter changes $\Delta\bm{\theta}$ and behavioral changes $\Delta p_{\bm{\theta}}$ is nonlinear, non-convex, and not characterizable in closed form. No transfer function exists. Without a plant model, no controller can compute corrective inputs.
    \item \textbf{Measurable error signal.} Feedback control requires $e(t) = r(t) - y(t)$, where $r(t)$ is a reference signal. The reference signal for ``general intelligence'' would require a complete specification of AGI---which loops back to the verification problem of Theorem~\ref{thm:godel}.
    \item \textbf{Stability guarantees.} Lyapunov stability and BIBO stability require structural properties (linearity, bounded gain, known dynamics) that LLM training dynamics lack. The self-training loop has no Lyapunov function guaranteeing convergence.
\end{enumerate}

A PID controller on a motor works because you can measure the angle, you know the transfer function, and the error signal is trivially defined. An LLM fails on all three. The system is unobservable, uncharacterizable, and has no well-defined reference signal. Self-regulation under these conditions is not an engineering problem to be solved---it is a category error.

\subsection{The Multi-Agent Correlation Problem}
\label{sec:multiagent}

\begin{objection}[Multiple Agents Checking Each Other]
What if we deploy multiple LLM agents in a pipeline---an ``architect'' that diagnoses problems and a ``programmer'' that implements fixes, or a ``generator'' and a ``critic''---so that they provide independent error correction for each other?
\end{objection}

\begin{response}
This does not work because the agents' errors are correlated. Multiple agents derived from the same (or similar) training distributions share:
\begin{itemize}
    \item The same training data $\implies$ the same blind spots,
    \item The same pretraining objective $\implies$ the same systematic biases,
    \item The same failure modes on edge cases and distribution-boundary inputs.
\end{itemize}

Let $\delta_{\mathrm{gen}}$ be the error introduced by the generator and $\delta_{\mathrm{crit}}$ be the error in the critic's evaluation. For independent correction, we need:
\begin{equation}
\mathrm{Corr}(\delta_{\mathrm{gen}}, \delta_{\mathrm{crit}}) \approx 0
\end{equation}
But since both agents are functions of the same training distribution $\mathcal{D}_{\mathrm{train}}$:
\begin{equation}
\mathrm{Corr}(\delta_{\mathrm{gen}}, \delta_{\mathrm{crit}}) \to 1 \quad \text{as overlap in training data} \to 1
\end{equation}

When the generator introduces a subtle bug, the critic has the same blind spot because it learned from the same distribution. The errors \emph{resonate}---they reinforce rather than cancel.

This generalizes to any number of agents: $n$ agents with correlated errors provide strictly less than $n$ independent checks. In the limit of identical training distributions, $n$~agents provide exactly one independent perspective, regardless of~$n$. The correlation is not in the architecture---it is in the training distribution. You cannot escape it by factoring the agent into components.

The SWE-CI benchmark~\cite{chen2025sweci} provides direct empirical evidence: its dual-agent architecture (Architect + Programmer) still exhibits regressions in over 75\% of tasks despite the agents checking each other (see Section~\ref{sec:sweci}). Kim et al.~\cite{kim2025correlated} evaluate over 350 LLMs and find 60\% error agreement when both models err, with error correlation \emph{increasing} for larger, more accurate models---demonstrating that scaling amplifies rather than alleviates the shared blind spots that prevent reliable mutual verification.

The only way to decorrelate errors is to introduce a signal from outside the distribution---which is human influence by Definition~\ref{def:human}.
\end{response}

\subsection{Why Noise Cannot Substitute for Human Influence}
\label{sec:noise}

A natural objection to the multi-agent correlation problem is: ``Add randomness to decorrelate the agents.'' This fails because of Theorem~\ref{thm:info_ceiling}.

White noise is statistically independent of the model's errors but carries zero mutual information with ground truth:
\begin{equation}
I(\text{noise};\, p_{\mathrm{true}}) = 0
\end{equation}
Adding noise to the model's output:
\begin{equation}
I(p_{\bm{\theta}_k} + \text{noise};\, p_{\mathrm{true}}) \leq I(p_{\bm{\theta}_k};\, p_{\mathrm{true}})
\end{equation}
by the data processing inequality. Information is \emph{lost}, not gained. The noise decorrelates the agents' errors in random directions, not toward correctness. The result is agents that disagree more often, but not agents that disagree \emph{usefully}.

What is needed is not independence alone---it is \emph{informed} independence. The external corrector must have mutual information with ground truth that is statistically independent of the model's errors. A human reviewer satisfies this: their knowledge originates from a completely different process (biological cognition, real-world embodied experience) and correlates with ground truth through a different channel.

In control theory terms (Section~\ref{sec:control}): adding noise to the sensor does not improve the state estimate. What is needed is a \emph{second sensor} that measures the actual system through a different physical principle. Human oversight is that second sensor.

\subsection{Empirical Corroboration: The SWE-CI Benchmark}
\label{sec:sweci}

Independent empirical evidence for the theoretical results of this paper is provided by SWE-CI (Software Engineering---Continuous Integration)~\cite{chen2025sweci}, the first repository-level benchmark evaluating LLM agent capabilities in maintaining codebases over iterative CI loops.

\subsubsection{Experimental Setup}
SWE-CI comprises 100 tasks drawn from 68 real-world repositories, with each task spanning an average of 233~days and 71~consecutive commits of authentic evolutionary history. The benchmark evaluates 18~models from 8~providers, consuming over 10~billion tokens. The iteration cycle---run tests, define requirements, modify code, repeat---instantiates precisely the autonomous self-improvement loop of Definition~\ref{def:selfimprove}, applied to code rather than model weights.

SWE-CI employs an Architect--Programmer dual-agent protocol: the Architect identifies functional gaps and issues requirements; the Programmer implements them. This is an explicit attempt at independent error correction through role decomposition.

\subsubsection{Key Result: Zero-Regression Rate}
The zero-regression rate measures the fraction of tasks in which the agent never regresses---that is, never breaks a previously passing test throughout the entire maintenance process. This metric directly corresponds to whether $\delta_k > 0$ (undetectable degradation) occurs in our error model (Theorem~\ref{thm:error_compound}).

\textbf{Result: Most models achieve a zero-regression rate below 0.25.} In over 75\% of tasks, the agent introduces at least one regression. Only two models in the Claude Opus series exceed a zero-regression rate of~0.5, meaning even the best models regress in roughly half of all tasks. This is a direct, large-scale measurement of $\delta_k > 0$.

\subsubsection{Mapping to Theoretical Results}

\begin{table}[h]
\centering
\begin{tabular}{@{}lp{9.5cm}@{}}
\toprule
\textbf{Theorem} & \textbf{SWE-CI Evidence} \\
\midrule
\ref{thm:error_compound} (Error Divergence) & Zero-regression rate ${<}\,0.25$ for most models: regressions occur in 75\%+ of tasks, directly measuring $\delta_k > 0$ \\[4pt]
\ref{thm:collapse} (Distributional Collapse) & Agent quality degrades over longer iteration sequences; EvoScore with $\gamma > 1$ penalizes models that accumulate technical debt \\[4pt]
Sec.~\ref{sec:multiagent} (Correlated Corrector) & Dual-agent architecture (Architect + Programmer) still fails at comparable rates---role decomposition does not provide independent error correction \\[4pt]
Sec.~\ref{sec:control} (Human Influence) & Model improvement across the evaluation period comes only from human-driven releases (new architectures, training data, RLHF), not from autonomous iteration \\
\bottomrule
\end{tabular}
\caption{Mapping between theoretical impossibility results and SWE-CI empirical observations.}
\label{tab:sweci}
\end{table}

\subsubsection{Within-Run Degradation vs.\ Across-Release Improvement}

SWE-CI's results reveal a striking asymmetry. \emph{Within} an autonomous run (the CI iteration loop), quality degrades: regressions accumulate, technical debt compounds, and the agent's ability to maintain the codebase deteriorates over successive iterations. \emph{Across} model releases (Claude Sonnet to Claude Opus, GPT-4o to GPT-o3, etc.), performance improves steadily.

This asymmetry is consistent with what the theory predicts. Autonomous iteration (no human influence) leads to error divergence (Theorem~\ref{thm:error_compound}) and distributional collapse (Theorem~\ref{thm:collapse}). Inter-release improvement reflects human-driven development: new architectures, curated training data, human-designed reward functions, and engineering effort. The intelligence driving progress across releases originates in human cognition, supporting the theorem target in Section~\ref{sec:human_target}.

The SWE-CI benchmark thus provides empirical corroboration, across 18~models and 100~tasks, for several failure modes targeted by this proof program~\cite{chen2025sweci,shumailov2024model}.

%% ============================================================
\section{Empirical Evidence: Support, Tension, and Falsification}
\label{sec:empirical}
%% ============================================================

A proof program that only collects confirming evidence is not a proof program. This section audits the 2023--2026 record in both directions: results consistent with the claims, results in tension with them, and the criteria under which the program should be considered refuted. Evidence is classified as \emph{theoretical} (a proof under stated assumptions), \emph{benchmark} (a measured evaluation), or \emph{mechanistic} (a controlled experiment isolating a cause), and is time-indexed by the architecture generation it was measured on---a distinction that matters, because several of the strongest props in earlier versions of this manuscript were measured on pre-RLVR models.

\subsection{Evidence consistent with the claims}

\begin{table}[h]
\centering
\small
\begin{tabular}{@{}p{2.6cm}p{5.1cm}p{6.4cm}@{}}
\toprule
\textbf{Claim} & \textbf{Result} & \textbf{Content and generation tested} \\
\midrule
Info.\ ceiling (Thm~\ref{thm:info_ceiling}) &
Yue et al.~2025~\cite{yue2025rlvr} &
RLVR-trained models beat base models at small $k$ but \emph{lose} at large $k$ on pass@$k$: the reasoning is sampled from, and bounded by, the base model. Benchmark; 2025 reasoning models. \\[3pt]
& Huang et al.~2024~\cite{huang2024sharpening} &
Self-improvement formalized as \emph{sharpening}---concentrating mass on outputs the model already ranks highly---with limits set by the base policy's coverage. Theoretical. \\[3pt]
& Gudibande et al.~2023~\cite{gudibande2023false} &
Imitation of a stronger model transfers style, not capability. Benchmark; pre-RLVR. \\[3pt]
& Zenil~2026~\cite{limitsself2025} &
Cross-entropy training does not implement Solomonoff induction; data processing and Levin non-growth hold; divergence for finite learners. Theoretical. \\
\midrule
Collapse (Thm~\ref{thm:collapse}) &
Shumailov et al.~2024~\cite{shumailov2024model}; Alemohammad et al.~\cite{alemohammad2023mad}; Dohmatob et al.~\cite{dohmatob2024tails,dohmatob2024strong}; Seddik et al.~\cite{seddik2024synthetic} &
Recursive training on replaced synthetic data degrades monotonically; tails vanish first; 1\% contamination suffices and scale amplifies rather than mitigates. Mechanistic + theoretical; 2023--2025. \\
\midrule
Verification limits (Thm~\ref{thm:godel}) &
Stechly et al.~2024~\cite{stechly2024verification}; Huang et al.~2023~\cite{huang2023selfcorrect} &
Intrinsic self-critique degrades reasoning accuracy; critique content is largely irrelevant. Mechanistic; pre-RLVR (see \S\ref{sec:timeindex}). \\[3pt]
& Chen et al.~2026~\cite{chen2026rsi} &
Survey of self-improvement systems finds a verification hierarchy---formal verifiers strongest, intrinsic self-assessment weakest---and that demonstrated improvement tracks it. \\
\midrule
Correlated correctors (\S\ref{sec:multiagent}) &
Kim et al.~2025~\cite{kim2025correlated}; Cemri et al.~2025~\cite{cemri2025mas} &
60\% error agreement across 350+ models, rising with capability; a 14-mode taxonomy of multi-agent failures from 1600+ annotated traces, with verification failures a top cluster. Benchmark. \\[3pt]
& Pan et al.~2026~\cite{pan2026coverrl} &
The ``consensus trap'': maximizing self-consistency collapses output diversity and confidently reinforces systematic errors that evade detection---the paper's predicted failure mode, named and measured. \\
\midrule
Non-improvement (Thm~\ref{thm:error_compound}) &
Chen et al.~2026~\cite{chen2025sweci} &
Zero-regression rate below 0.25 for most of 18 models over 100 long-horizon maintenance tasks (\S\ref{sec:sweci}). Benchmark; 2026 agents. \\
\midrule
In-context self-modification (\S\ref{sec:context}) &
Zhang et al.~2026~\cite{zhang2026introspection} &
Sustained recursive self-improvement requires introspection; current LLMs achieve only quasi-introspection, blocked by incomplete self-access and the feedforward architecture. Theoretical. \\
\bottomrule
\end{tabular}
\caption{Results consistent with the claims of this manuscript.}
\label{tab:support}
\end{table}

\subsection{Evidence in tension with the claims}

The following results either contradict a claim made in earlier versions of this manuscript, or contradict a natural strengthening of a claim made here. Each is listed with the specific hypothesis it breaks.

\begin{table}[h]
\centering
\small
\begin{tabular}{@{}p{4.4cm}p{4.6cm}p{5.1cm}@{}}
\toprule
\textbf{Result} & \textbf{What it shows} & \textbf{Hypothesis it breaks} \\
\midrule
Gerstgrasser et al.~2024~\cite{gerstgrasser2024accumulate} &
Accumulating synthetic data alongside real data gives test error a finite bound independent of iteration count. &
Replacement regime (Def.~\ref{def:selfimprove}); refutes universal error divergence (\S\ref{sec:regimes}). \\[3pt]
Feng et al.~2024~\cite{feng2024verification}; Yi et al.~2026~\cite{yi2026escaping} &
Even imperfect verifiers prevent collapse; verifier-guided retraining converges to the \emph{verifier's} knowledge center, with gains that plateau or reverse as the verifier is imperfect. &
Collapse as inevitable; supplies the quantitative form of the grounding claim (\S\ref{sec:grounding_discussion}). \\[3pt]
Kang et al.~2025~\cite{kang2025demystifying} &
Rephrased synthetic data mixed at $\sim$1/3 shows no degradation at foreseeable scales and speeds up training 5--10$\times$; textbook-style generated data does collapse. &
Collapse claims stated without regard to synthesis method and mixing ratio. \\[3pt]
Zhang et al.~2025 (DGM)~\cite{zhang2025dgm} &
An agent rewriting its own code raises SWE-bench from 20.0\% to 50.0\% and Polyglot from 14.2\% to 30.7\%. &
Nothing formally---the loop is scored by human-written test suites---but it refutes casual readings of ``LLMs cannot self-improve''. \\[3pt]
Zhang et al.~2026 (Self-Harness)~\cite{zhang2026selfharness} &
Agents mine their own traces, propose harness changes, and validate by regression testing, improving held-in \emph{and} held-out pass rates in all nine model--benchmark pairs, up to 132\% relative. &
Same: grounded by benchmark suites, but the human contribution per iteration is close to zero. \\[3pt]
Zhao et al.~2025 (AZR)~\cite{zhao2025azr}; Huang et al.~2025 (R-Zero)~\cite{huang2025rzero} &
AZR proposes and solves its own coding tasks, validated by a code executor, and outperforms zero-setting models trained on tens of thousands of curated examples. R-Zero co-evolves a challenger and a solver with \emph{no} external verifier---pseudo-labels come from majority voting over the solver's own samples---for gains of $+6.5$ (math) and $+7.5$ (general reasoning) on Qwen3-4B-Base. &
R-Zero in particular has \emph{no} external correctness signal, which is the closest published approach to the sealed loop. \\[3pt]
Liu et al.~2025 (ProRL)~\cite{liu2025prorl} &
Prolonged RL finds solutions inaccessible to the base model at any $k$, contradicting Yue et al.~\cite{yue2025rlvr}. &
The strong reading of the information ceiling as a capability ceiling. \\[3pt]
Chen et al.~2026~\cite{chen2026selfcorrection} &
Relabeling an error's apparent \emph{source} (internal thought vs.\ external message), with content held fixed, raises explicit correction rates by 23--93 points in 10 of 12 settings. &
The claim that models cannot detect their own errors: much of the 2023 self-correction failure is a presentation artifact, not a capability limit. \\
\bottomrule
\end{tabular}
\caption{Results in tension with the claims, and the specific hypothesis each one breaks.}
\label{tab:tension}
\end{table}

\subsection{The grounding audit}
\label{sec:grounding_audit}

Table~\ref{tab:tension} invites an obvious objection: if every successful system is reclassified as ``externally grounded'', the thesis is unfalsifiable. The defense has to be that the grounding source is identifiable \emph{in advance} and independently of the outcome. Table~\ref{tab:audit} performs that audit for the strongest 2025--26 systems; in each case the source is a concrete artifact that existed before the loop started.

\begin{table}[h]
\centering
\small
\begin{tabular}{@{}p{3.2cm}p{5.4cm}p{5.5cm}@{}}
\toprule
\textbf{System} & \textbf{Grounding source (named in advance)} & \textbf{Status under Def.~\ref{def:grounding}} \\
\midrule
DGM~\cite{zhang2025dgm} & SWE-bench / Polyglot task suites with human-written tests & Grounded; human origin \\
AlphaEvolve~\cite{novikov2025alphaevolve} & Per-problem automated evaluators; correctness of the discovered algorithm is machine-checkable & Grounded; human-designed evaluator, non-human verification \\
SEAL~\cite{zweiger2025seal} & Downstream task performance on held-out evaluation data as RL reward & Grounded; human origin \\
Absolute Zero~\cite{zhao2025azr} & Code executor (deterministic ground truth for program behaviour) & Grounded; artifact origin, not per-task human labels \\
Self-Harness~\cite{zhang2026selfharness} & Regression tests against benchmark suites & Grounded; human origin, near-zero per-iteration cost \\
R-Zero~\cite{huang2025rzero} & \emph{None}: majority-vote pseudo-labels from the model itself & Ungrounded---the sealed case, and the one to watch \\
CoVerRL~\cite{pan2026coverrl} & \emph{None}: self-verification, bootstrapped & Ungrounded; documents the consensus trap that the sealed case predicts \\
\bottomrule
\end{tabular}
\caption{Where the information comes from in demonstrated self-improvement systems.}
\label{tab:audit}
\end{table}

The pattern is consistent and is the strongest empirical statement this program can honestly make: \emph{the magnitude of demonstrated self-improvement tracks the strength of the external signal available}, from machine-checkable execution at the top to self-consistency at the bottom, exactly as Chen et al.~\cite{chen2026rsi} report for the field as a whole. The two ungrounded systems are also the two that report the smallest gains and the failure mode this manuscript predicts---diversity collapse with confident reinforcement of systematic error---which R-Zero mitigates with a self-generated curriculum and CoVerRL mitigates by bootstrapping a verifier, in both cases without an external correctness signal and in both cases with gains in the single-digit percentage range on benchmarks built by humans.

\subsection{Time-indexing the older evidence}
\label{sec:timeindex}

Three empirical props used in earlier versions require qualification.

\textbf{Intrinsic self-correction.} Huang et al.~\cite{huang2023selfcorrect} and Stechly et al.~\cite{stechly2024verification} measured GPT-3.5/GPT-4-era models without verifier training. Chen et al.~\cite{chen2026selfcorrection} show that a large part of the observed failure is an artifact of how errors are labeled in the chat template: relabeling the same erroneous content as externally sourced raises explicit correction rates by up to 93 points. The defensible current claim is narrower: \emph{unaided self-verification is unreliable and its reliability is sensitive to presentation}, not that models cannot detect their own errors.

\textbf{The reversal curse and compositional collapse.} Berglund et al.~\cite{berglund2023reversal} and Dziri et al.~\cite{dziri2023faith} were measured on 2023 models. Both effects are substantially attenuated in verifier-trained reasoning models. They remain evidence that autoregressive pattern completion does not automatically yield systematic reasoning; they are no longer evidence about the capability of current systems.

\textbf{NP-hard benchmarks.} Fan et al.~\cite{fan2023nphardeval} likewise predates reasoning models. The manuscript already declines to treat these failures as evidence for $\mathrm{P} \neq \mathrm{NP}$ (Remark~\ref{rem:pnp_empirical}); they are now also weak evidence about current capability.

\subsection{What would falsify this program}
\label{sec:falsification}

The following are stated in advance and would each require a substantive revision or retraction.

\begin{enumerate}
    \item \textbf{Sealed-loop capability growth.} A loop satisfying Definition~\ref{def:selfimprove}---no external verifier, no new data, replacement dynamics---showing monotone improvement over at least five iterations on a held-out benchmark scored externally, with the improvement exceeding the base model's pass@$k$ coverage at the compute used. This would refute the practical content of Corollary~\ref{cor:accuracy_ceiling}. R-Zero~\cite{huang2025rzero} is the nearest existing attempt---majority-vote pseudo-labels, no external verifier---and we are not aware of a published pass@$k$ comparison against its base model at matched compute. Running that comparison is a concrete, cheap experiment that would bear directly on this criterion.
    \item \textbf{Non-saturating grounding.} A demonstration that capability keeps rising past the competence of the grounding source---contradicting the convergence-to-verifier-knowledge-center result of Yi et al.~\cite{yi2026escaping}---would refute the bits-of-grounding conjecture of \S\ref{sec:grounding_discussion}.
    \item \textbf{Broken channel model.} A construction in which the training pipeline is \emph{not} a channel in the sense of Definition~\ref{def:joint}---for instance, where the pseudo-random stream used in sampling is correlated with $\Theta$---would break the data-processing hypothesis on which everything rests. This is not a curiosity: a model whose sampling seed is derived from real-world data is grounded, quietly.
    \item \textbf{Target reached from below.} An autonomous loop reaching the criterion of Definition~\ref{def:agi} on a fixed, pre-registered family from a starting model with $I_0$ demonstrably below the requirement of Corollary~\ref{cor:accuracy_ceiling}.
\end{enumerate}

What would \emph{not} count: gains from systems whose loop consults human-built verifiers, test suites, compilers, proof checkers, or benchmark scores; gains that disappear on held-out task families; and gains measured only by the improving system itself.

\subsection{Benchmark trends, and the limits of information accounting}

Two trend series bear on the manuscript's claims without settling them. ARC-AGI-2~\cite{chollet2025arcagi2}, on which every frontier model scored 0\% at its March 2025 launch, saw a top Kaggle score of 24\% on the private evaluation set within the first competition cycle~\cite{arcprize2025}, with higher verified scores reported for commercial systems. METR's task-length measurements~\cite{kwa2025horizons} put the 50\% time horizon of frontier models at roughly 50 minutes in early 2025 and, in the January 2026 revision of the task suite~\cite{metr2026horizon}, at 320 minutes for the strongest measured model, with doubling times of 131 days when fitted from 2023 and 89 days when fitted from 2024.

Both series measure \emph{across-release} progress, which is human-driven: new architectures, new curated corpora, new human-designed reward schemes. They are therefore consistent with the within-run/across-release asymmetry documented in \S\ref{sec:sweci}, and they are not evidence about sealed autonomous loops. They do, however, carry a warning for this manuscript. Capability can jump non-linearly with no corresponding jump in information, because accessibility of latent information changes---the distinction Corollary~\ref{cor:pnp_info} draws. Information accounting bounds what a sealed loop can know; it is a poor instrument for forecasting what a grounded system will be able to do, and no forecast of that kind should be read into these results.

%% ============================================================
%% Bibliography
%% ============================================================

\begin{thebibliography}{99}

\bibitem{shannon1948}
C.~E. Shannon.
\newblock A mathematical theory of communication.
\newblock \emph{Bell System Technical Journal}, 27(3):379--423, 1948.

\bibitem{cover2006}
T.~M. Cover and J.~A. Thomas.
\newblock \emph{Elements of Information Theory}.
\newblock Wiley-Interscience, 2nd edition, 2006.

\bibitem{godel1931}
K.~G\"{o}del.
\newblock \"{U}ber formal unentscheidbare {S}\"{a}tze der {P}rincipia {M}athematica und verwandter {S}ysteme {I}.
\newblock \emph{Monatshefte f\"{u}r Mathematik und Physik}, 38:173--198, 1931.

\bibitem{rice1953}
H.~G. Rice.
\newblock Classes of recursively enumerable sets and their decision problems.
\newblock \emph{Transactions of the American Mathematical Society}, 74(2):358--366, 1953.

\bibitem{shumailov2024model}
I.~Shumailov, Z.~Shumaylov, Y.~Zhao, N.~Papernot, R.~Anderson, and Y.~Gal.
\newblock {AI} models collapse when trained on recursively generated data.
\newblock \emph{Nature}, 631:755--759, 2024.

\bibitem{bender2021stochastic}
E.~M. Bender, T.~Gebru, A.~McMillan-Major, and S.~Shmitchell.
\newblock On the dangers of stochastic parrots: Can language models be too big?
\newblock In \emph{Proceedings of FAccT}, pages 610--623, 2021.

\bibitem{chen2025sweci}
J.~Chen, X.~Xu, H.~Wei, C.~Chen, and B.~Zhao.
\newblock {SWE-CI}: Evaluating agent capabilities in maintaining codebases via continuous integration.
\newblock \emph{arXiv preprint arXiv:2603.03823}, 2026.

\bibitem{silver2017alphazero}
D.~Silver, T.~Hubert, J.~Schrittwieser, I.~Antonoglou, M.~Lai, A.~Guez, M.~Lanctot, L.~Sifre, D.~Kumaran, T.~Graepel, T.~Lillicrap, K.~Simonyan, and D.~Hassabis.
\newblock Mastering chess and shogi by self-play with a general reinforcement learning algorithm.
\newblock \emph{arXiv preprint arXiv:1712.01815}, 2017.

\bibitem{vaswani2017attention}
A.~Vaswani, N.~Shazeer, N.~Parmar, J.~Uszkoreit, L.~Jones, A.~N. Gomez, \L.~Kaiser, and I.~Polosukhin.
\newblock Attention is all you need.
\newblock In \emph{Advances in Neural Information Processing Systems}, volume~30, 2017.

\bibitem{penrose1994}
R.~Penrose.
\newblock \emph{Shadows of the Mind: A Search for the Missing Science of Consciousness}.
\newblock Oxford University Press, 1994.

\bibitem{bostrom2014}
N.~Bostrom.
\newblock \emph{Superintelligence: Paths, Dangers, Strategies}.
\newblock Oxford University Press, 2014.

\bibitem{ouyang2022rlhf}
L.~Ouyang, J.~Wu, X.~Jiang, D.~Almeida, C.~Wainwright, P.~Mishkin, C.~Zhang, S.~Agarwal, K.~Slama, A.~Ray, et~al.
\newblock Training language models to follow instructions with human feedback.
\newblock In \emph{Advances in Neural Information Processing Systems}, volume~35, 2022.

\bibitem{kolmogorov1965}
A.~N. Kolmogorov.
\newblock Three approaches to the quantitative definition of information.
\newblock \emph{Problems of Information Transmission}, 1(1):1--7, 1965.

%% ===== New references: Empirical evidence =====

\bibitem{gudibande2023false}
A.~Gudibande, E.~Wallace, C.~Snell, X.~Geng, H.~Liu, P.~Abbeel, S.~Levine, and D.~Song.
\newblock The false promise of imitating proprietary {LLMs}.
\newblock In \emph{International Conference on Learning Representations}, 2024.
\newblock arXiv:2305.15717.

\bibitem{seddik2024synthetic}
M.~E.~A. Seddik, S.-W.~Chen, S.~Hayou, P.~Youssef, and M.~Debbah.
\newblock How bad is training on synthetic data? {A} statistical analysis of language model collapse.
\newblock \emph{arXiv preprint arXiv:2404.05090}, 2024.

\bibitem{alemohammad2023mad}
S.~Alemohammad, J.~Casco-Rodriguez, L.~Luzi, A.~I. Humayun, H.~Babaei, D.~LeJeune, A.~Siahkoohi, and R.~G. Baraniuk.
\newblock Self-consuming generative models go {MAD}.
\newblock In \emph{International Conference on Learning Representations}, 2024.
\newblock arXiv:2307.01850.

\bibitem{dohmatob2024tails}
E.~Dohmatob, Y.~Feng, P.~Yang, F.~Charton, and J.~Kempe.
\newblock A tale of tails: Model collapse as a change of scaling laws.
\newblock In \emph{International Conference on Machine Learning}, 2024.
\newblock arXiv:2402.07043.

\bibitem{dohmatob2024strong}
E.~Dohmatob, Y.~Feng, A.~Subramonian, and J.~Kempe.
\newblock Strong model collapse.
\newblock In \emph{International Conference on Learning Representations}, 2025.
\newblock arXiv:2410.04840.

\bibitem{huang2023selfcorrect}
J.~Huang, X.~Chen, S.~Mishra, H.~S. Zheng, A.~W. Yu, X.~Song, and D.~Zhou.
\newblock Large language models cannot self-correct reasoning yet.
\newblock In \emph{International Conference on Learning Representations}, 2024.
\newblock arXiv:2310.01798.

\bibitem{zheng2023judge}
L.~Zheng, W.-L.~Chiang, Y.~Sheng, S.~Zhuang, Z.~Wu, Y.~Zhuang, Z.~Lin, Z.~Li, D.~Li, E.~P.~Xing, H.~Zhang, J.~E.~Gonzalez, and I.~Stoica.
\newblock Judging {LLM}-as-a-judge with {MT-Bench} and {Chatbot Arena}.
\newblock In \emph{Advances in Neural Information Processing Systems}, volume~36, 2023.

\bibitem{stechly2024verification}
K.~Stechly, K.~Valmeekam, and S.~Kambhampati.
\newblock On the self-verification limitations of large language models on reasoning and planning tasks.
\newblock \emph{arXiv preprint arXiv:2402.08115}, 2024.

\bibitem{babcock2024superradiance}
N.~S.~Babcock, G.~Montes-Cabrera, K.~E.~Oberhofer, M.~Chergui, G.~L.~Celardo, and P.~Kurian.
\newblock Ultraviolet superradiance from mega-networks of tryptophan in biological architectures.
\newblock \emph{The Journal of Physical Chemistry B}, 128(17):4035--4046, 2024.

\bibitem{kalra2023exciton}
A.~P.~Kalra, A.~Benny, S.~M.~Travis, E.~A.~Zizzi, A.~Morales-Sanchez, D.~G.~Oblinsky, T.~J.~A.~Craddock, S.~R.~Hameroff, M.~B.~MacIver, J.~A.~Tuszynski, S.~Petry, R.~Penrose, and G.~D.~Scholes.
\newblock Electronic energy migration in microtubules.
\newblock \emph{ACS Central Science}, 9(3):352--361, 2023.

\bibitem{wiest2024epothilone}
M.~C.~Wiest et~al.
\newblock Microtubule-stabilizer epothilone {B} delays anesthetic-induced unconsciousness in rats.
\newblock \emph{eNeuro}, 11(8), 2024.

\bibitem{li2018xenon}
N.~Li, D.~Lu, L.~Yang, H.~Tao, Y.~Xu, C.~Wang, L.~Fu, H.~Liu, Y.~Chummum, and S.~Zhang.
\newblock Nuclear spin attenuates the anesthetic potency of xenon isotopes in mice.
\newblock \emph{Anesthesiology}, 129(2):271--277, 2018.

\bibitem{straub2025posner}
J.~S.~Straub, M.~L.~Patel, M.~S.~Nowotarski, L.~Rao, M.~E.~Turiansky, M.~P.~A.~Fisher, and M.~E.~Helgeson.
\newblock Evidence for a possible quantum effect on calcium phosphate formation.
\newblock \emph{Proceedings of the National Academy of Sciences}, 122(10):e2423211122, 2025.

\bibitem{hameroff2014orchor}
S.~Hameroff and R.~Penrose.
\newblock Consciousness in the universe: A review of the `{O}rch {OR}' theory.
\newblock \emph{Physics of Life Reviews}, 11:39--78, 2014.

\bibitem{furlanello2018born}
T.~Furlanello, Z.~Lipton, M.~Tschannen, L.~Itti, and A.~Anandkumar.
\newblock Born again neural networks.
\newblock In \emph{International Conference on Machine Learning}, 2018.

\bibitem{mobahi2020selfdistillation}
H.~Mobahi, M.~Farajtabar, and P.~L.~Bartlett.
\newblock Self-distillation amplifies regularization in {Hilbert} space.
\newblock In \emph{Advances in Neural Information Processing Systems}, volume~33, 2020.

\bibitem{dziri2023faith}
N.~Dziri, X.~Lu, M.~Sclar, X.~L.~Li, L.~Jiang, B.~Y.~Lin, P.~West, C.~Bhagavatula, R.~Le~Bras, J.~D.~Hwang, S.~Sanyal, S.~Welleck, X.~Ren, A.~Ettinger, Z.~Harchaoui, and Y.~Choi.
\newblock Faith and fate: Limits of transformers on compositionality.
\newblock In \emph{Advances in Neural Information Processing Systems}, volume~36, 2023.

\bibitem{berglund2023reversal}
L.~Berglund, M.~Tong, M.~Kaufmann, M.~Balesni, A.~C.~Stickland, T.~Korbak, and O.~Evans.
\newblock The reversal curse: {LLMs} trained on ``{A} is {B}'' fail to learn ``{B} is {A}''.
\newblock In \emph{International Conference on Learning Representations}, 2024.

\bibitem{fan2023nphardeval}
L.~Fan, W.~Hua, L.~Li, H.~Ling, and Y.~Zhang.
\newblock {NPHardEval}: Dynamic benchmark on reasoning ability of large language models via complexity classes.
\newblock \emph{arXiv preprint arXiv:2312.14890}, 2023.

\bibitem{liu2023lost}
N.~F.~Liu, K.~Lin, J.~Hewitt, A.~Paranjape, M.~Bevilacqua, F.~Petroni, and P.~Liang.
\newblock Lost in the middle: How language models use long contexts.
\newblock \emph{Transactions of the Association for Computational Linguistics}, 2024.

\bibitem{kaplan2020scaling}
J.~Kaplan, S.~McCandlish, T.~Henighan, T.~B.~Brown, B.~Chess, R.~Child, S.~Gray, A.~Radford, J.~Wu, and D.~Amodei.
\newblock Scaling laws for neural language models.
\newblock \emph{arXiv preprint arXiv:2001.08361}, 2020.

\bibitem{hoffmann2022chinchilla}
J.~Hoffmann, S.~Borgeaud, A.~Mensch, E.~Buchatskaya, T.~Cai, E.~Rutherford, D.~de~Las~Casas, L.~A.~Hendricks, J.~Welbl, A.~Clark, et~al.
\newblock Training compute-optimal large language models.
\newblock \emph{arXiv preprint arXiv:2203.15556}, 2022.

\bibitem{gao2022overoptimization}
L.~Gao, J.~Schulman, and J.~Hilton.
\newblock Scaling laws for reward model overoptimization.
\newblock In \emph{International Conference on Machine Learning}, 2023.

\bibitem{limitsself2025}
H.~Zenil.
\newblock On {S}olomonoff induction in large language models and the limits of self-improving: The singularity is not near without symbolic model synthesis.
\newblock \emph{arXiv preprint arXiv:2601.05280}, 2026. Version~3, August 2026.

\bibitem{gerstgrasser2024accumulate}
M.~Gerstgrasser, R.~Schaeffer, A.~Dey, R.~Rafailov, H.~Sleight, J.~Hughes, T.~Korbak, R.~Agrawal, D.~Pai, A.~Gromov, D.~A.~Roberts, D.~Yang, D.~L.~Donoho, and S.~Koyejo.
\newblock Is model collapse inevitable? {B}reaking the curse of recursion by accumulating real and synthetic data.
\newblock \emph{arXiv preprint arXiv:2404.01413}, 2024.

\bibitem{feng2024verification}
Y.~Feng, E.~Dohmatob, P.~Yang, F.~Charton, and J.~Kempe.
\newblock Beyond model collapse: Scaling up with synthesized data requires verification.
\newblock \emph{arXiv preprint arXiv:2406.07515}, 2024.

\bibitem{yi2026escaping}
B.~Yi, Q.~Liu, Y.~Cheng, and H.~Xu.
\newblock Escaping model collapse via synthetic data verification: Near-term improvements and long-term convergence.
\newblock \emph{arXiv preprint arXiv:2510.16657}, 2025. Revised 2026.

\bibitem{kang2025demystifying}
F.~Kang, N.~Ardalani, M.~Kuchnik, Y.~Emad, M.~Elhoushi, S.~Sengupta, S.-W.~Li, R.~Raghavendra, R.~Jia, and C.-J.~Wu.
\newblock Demystifying synthetic data in {LLM} pre-training: A systematic study of scaling laws, benefits, and pitfalls.
\newblock In \emph{Proceedings of EMNLP}, 2025.
\newblock arXiv:2510.01631.

\bibitem{yue2025rlvr}
Y.~Yue, Z.~Chen, R.~Lu, A.~Zhao, Z.~Wang, S.~Song, and G.~Huang.
\newblock Does reinforcement learning really incentivize reasoning capacity in {LLM}s beyond the base model?
\newblock In \emph{Advances in Neural Information Processing Systems}, 2025.
\newblock arXiv:2504.13837.

\bibitem{liu2025prorl}
M.~Liu, S.~Diao, X.~Lu, J.~Hu, X.~Dong, Y.~Choi, J.~Kautz, and Y.~Dong.
\newblock {ProRL}: Prolonged reinforcement learning expands reasoning boundaries in large language models.
\newblock \emph{arXiv preprint arXiv:2505.24864}, 2025.

\bibitem{huang2024sharpening}
A.~Huang, A.~Block, D.~J.~Foster, D.~Rohatgi, C.~Zhang, M.~Simchowitz, J.~T.~Ash, and A.~Krishnamurthy.
\newblock Self-improvement in language models: The sharpening mechanism.
\newblock \emph{arXiv preprint arXiv:2412.01951}, 2024.

\bibitem{song2024gap}
Y.~Song, H.~Zhang, C.~Eisenach, S.~Kakade, D.~Foster, and U.~Ghai.
\newblock Mind the gap: Examining the self-improvement capabilities of large language models.
\newblock In \emph{International Conference on Learning Representations}, 2025.
\newblock arXiv:2412.02674.

\bibitem{zhang2025dgm}
J.~Zhang, S.~Hu, C.~Lu, R.~Lange, and J.~Clune.
\newblock Darwin {G}\"{o}del machine: Open-ended evolution of self-improving agents.
\newblock In \emph{International Conference on Learning Representations}, 2026.
\newblock arXiv:2505.22954.

\bibitem{zweiger2025seal}
A.~Zweiger, J.~Pari, H.~Guo, E.~Aky\"{u}rek, Y.~Kim, and P.~Agrawal.
\newblock Self-adapting language models.
\newblock \emph{arXiv preprint arXiv:2506.10943}, 2025.

\bibitem{zhao2025azr}
A.~Zhao, Y.~Wu, Y.~Yue, T.~Wu, Q.~Xu, Y.~Yue, M.~Lin, S.~Wang, Q.~Wu, Z.~Zheng, and G.~Huang.
\newblock Absolute zero: Reinforced self-play reasoning with zero data.
\newblock In \emph{Advances in Neural Information Processing Systems}, 2025.
\newblock arXiv:2505.03335.

\bibitem{huang2025rzero}
C.~Huang, W.~Yu, X.~Wang, H.~Zhang, Z.~Li, R.~Li, J.~Huang, H.~Mi, and D.~Yu.
\newblock {R-Zero}: Self-evolving reasoning {LLM} from zero data.
\newblock \emph{arXiv preprint arXiv:2508.05004}, 2025.

\bibitem{zhang2026selfharness}
H.~Zhang, S.~Zhang, K.~Li, C.~Zhang, Y.~Chen, Y.~Zhang, L.~Bai, and S.~Hu.
\newblock Self-harness: Harnesses that improve themselves.
\newblock \emph{arXiv preprint arXiv:2606.09498}, 2026.

\bibitem{novikov2025alphaevolve}
A.~Novikov, N.~V\~{u}, M.~Eisenberger, E.~Dupont, P.-S.~Huang, A.~Z.~Wagner, et~al.
\newblock {AlphaEvolve}: A coding agent for scientific and algorithmic discovery.
\newblock \emph{arXiv preprint arXiv:2506.13131}, 2025.

\bibitem{chen2026selfcorrection}
K.-Y.~Chen, F.-Y.~Su, S.-Y.~Lin, B.~Li, and J.-H.~Chiang.
\newblock The self-correction illusion: Role relabeling gates explicit error flagging in large language models.
\newblock \emph{arXiv preprint arXiv:2606.05976}, 2026.

\bibitem{pan2026coverrl}
T.~Pan, Y.~Yan, Z.~Wang, R.~Zhang, G.~Hou, W.~Zhang, W.~Lu, J.~Xiao, and Y.~Shen.
\newblock {CoVerRL}: Breaking the consensus trap in label-free reasoning via generator--verifier co-evolution.
\newblock \emph{arXiv preprint arXiv:2603.17775}, 2026.

\bibitem{zhang2026introspection}
J.~Zhang, B.~Yuan, and Q.~Zhang.
\newblock Self-reference in large language models: The introspection threshold for recursive self-improvement.
\newblock \emph{arXiv preprint arXiv:2607.04277}, 2026.

\bibitem{chen2026rsi}
M.~Chen, L.~Wang, and B.~Qu.
\newblock Recursive self-improvement in {AI}: From bounded self-refinement to autonomous research loops.
\newblock \emph{arXiv preprint arXiv:2607.07663}, 2026.

\bibitem{cemri2025mas}
M.~Cemri, M.~Z.~Pan, S.~Yang, L.~A.~Agrawal, B.~Chopra, R.~Tiwari, K.~Keutzer, A.~Parameswaran, D.~Klein, K.~Ramchandran, M.~Zaharia, J.~E.~Gonzalez, and I.~Stoica.
\newblock Why do multi-agent {LLM} systems fail?
\newblock \emph{arXiv preprint arXiv:2503.13657}, 2025.

\bibitem{chollet2025arcagi2}
F.~Chollet, M.~Knoop, G.~Kamradt, B.~Landers, and H.~Pinkard.
\newblock {ARC-AGI-2}: A new challenge for frontier {AI} reasoning systems.
\newblock \emph{arXiv preprint arXiv:2505.11831}, 2025.

\bibitem{arcprize2025}
F.~Chollet, M.~Knoop, G.~Kamradt, and B.~Landers.
\newblock {ARC} {P}rize 2025: Technical report.
\newblock \emph{arXiv preprint arXiv:2601.10904}, 2026.

\bibitem{kwa2025horizons}
T.~Kwa, B.~West, J.~Becker, A.~Deng, K.~Garcia, et~al. (METR).
\newblock Measuring {AI} ability to complete long software tasks.
\newblock \emph{arXiv preprint arXiv:2503.14499}, 2025.

\bibitem{metr2026horizon}
METR.
\newblock Time horizon 1.1.
\newblock Technical report, 29 January 2026.
\newblock \url{https://metr.org/blog/2026-1-29-time-horizon-1-1/}.

\bibitem{kim2025correlated}
E.~Kim et~al.
\newblock Correlated errors in large language models.
\newblock \emph{arXiv preprint arXiv:2506.07962}, 2025.

\end{thebibliography}

\end{document}
